\documentclass[11pt]{article}

\usepackage[english]{babel}
\usepackage[top=1in,bottom=1.2in,left=1in,right=1in]{geometry}
\usepackage{amssymb,amsmath,amsthm,mathtools}
\usepackage{mathrsfs}
\usepackage{times}
\usepackage{graphicx}
\usepackage{placeins}
\usepackage{hyperref}

\hypersetup{
  colorlinks=true,
  linkcolor=blue,
  citecolor=red,
  urlcolor=blue,
  pdfauthor={}
}

\newtheorem{theorem}{Theorem}[section]
\newtheorem{proposition}[theorem]{Proposition}
\newtheorem{lemma}[theorem]{Lemma}
\newtheorem{corollary}[theorem]{Corollary}
\theoremstyle{definition}
\newtheorem{remark}[theorem]{Remark}

\numberwithin{equation}{section}

\newcommand{\E}{\mathbb{E}}
\newcommand{\R}{\mathbb{R}}
\newcommand{\cN}{\mathcal{N}}
\newcommand{\wh}{\widehat}
\newcommand{\eps}{\varepsilon}
\newcommand{\op}{\mathrm{op}}
\newcommand{\Id}{\mathrm{Id}}
\DeclareMathOperator{\Var}{Var}
\DeclareMathOperator{\Cov}{Cov}
\DeclareMathOperator{\supp}{supp}
\DeclareMathOperator{\conv}{conv}

\title{Support Divergence of Mixture NPMLEs Under Finite Mixtures}
\author{}
\date{}

\begin{document}

\maketitle

\begin{abstract}
Motivated by distribution testing and empirical Bayes calibration, we ask how sparse a mixture NPMLE can be when the true mixing distribution has fixed finite support. For Gaussian location mixtures in every fixed dimension \(d\), every maximizer has at least \(c(\log n)^{d/2}/\log\log n\) atoms with high probability. For Poisson mixtures, the lower bound is \(c\sqrt{\log n}/(\log\log n)^{3/2}\). Thus the NPMLE can have diverging support even under a fixed finite mixture.
\end{abstract}

\medskip
\noindent\textbf{Note.} This draft was produced through an iterative process in which Mark Sellke repeatedly asked GPT-5.6 in OpenAI Codex to prove, refine, and write up further extensions.
\medskip

\setcounter{tocdepth}{2}
\tableofcontents

\section{Introduction}

Mixture models describe heterogeneous populations through a probability
distribution \(G\) over latent parameters. Writing
\(p_G(x)=\int p_\theta(x)\,dG(\theta)\), each atom of \(G\) corresponds to a
component of the mixture. When the number of components is unknown, the
nonparametric maximum likelihood estimator (NPMLE) maximizes the likelihood
over all probability measures \(G\), rather than fixing a support size in
advance. This construction traces to Kiefer--Wolfowitz
\cite{kiefer1956consistency} and has closely related empirical-Bayes roots in
Robbins \cite{robbins1950generalization}: there \(G\) is an unknown prior,
and the fitted mixing distribution is used to estimate posterior quantities
or Bayes rules. Yet the mixture density, the induced decision rule, and the
support of the mixing distribution are different statistical objects.

We ask whether the atom count of the NPMLE adapts to a finitely supported
truth. Suppose the data law is \(p_{G_0}\) for a fixed finite mixture
\(G_0\). For Gaussian location and Poisson mixtures, we show that the support
of every NPMLE diverges with the sample size. Thus accurate density or
empirical-Bayes estimation can coexist with a fitted atom count that does not
estimate the number of latent components. Approximate, thresholded, and
penalized procedures define different estimators whose atom counts may
depend on stopping or tuning rules.

The question was posed explicitly for one-dimensional Gaussian location
mixtures by Polyanskiy and Wu \cite{polyanskiy2020self}. They proved the
self-regularization bound \(O_{\mathbb P}(\log n)\) for subgaussian data and
asked in \cite[Section~5.6]{polyanskiy2020self} whether this order is attained
even under a single Gaussian component. Lindsay's convex-geometric theory
gives a maximizer with at most \(n\) atoms \cite{lindsay1983geometryA}, and
the Lindsay--Roeder criterion makes the NPMLE unique in the one-dimensional
models considered here \cite{lindsay1993uniqueness}. Multivariate Gaussian
NPMLEs can be nonunique \cite{soloff2024multivariate}, so our Gaussian lower
bound is uniform over all maximizers. We extend the single-Gaussian question
to every fixed dimension and every fixed finite Gaussian location mixture,
and prove an analogous result for finite Poisson mixtures with positive
largest mean. Our companion papers determine how much variance inflation is
needed to collapse the Gaussian fit to one atom in dimension one and in
dimensions at least two
\cite{varianceD1Companion,varianceMultivariateCompanion}.

\subsection*{Preliminaries and notation}

All asymptotic notation is as $n\to\infty$, unless another limit is
specified. We write $\mathbb P$ and $\mathbb E$ for probability and
expectation under the distribution currently under discussion. For a
deterministic positive sequence $a_n$, the relation
$X_n=O_{\mathbb P}(a_n)$ means that $|X_n|/a_n$ is tight, while
$X_n=o_{\mathbb P}(a_n)$ means that $X_n/a_n\to0$ in probability.
Deterministic $O(\cdot)$ and $o(\cdot)$ have their usual meanings. We write
$a_n\sim b_n$ if $a_n/b_n\to1$, $a_n\ll b_n$ if $a_n/b_n\to0$, and
$a_n\lesssim b_n$ if $a_n\le Cb_n$ for a constant independent of $n$.
The reverse relations are denoted by $\gg$ and $\gtrsim$, and
$a_n\asymp b_n$ means that both one-sided bounds hold. Subscripts, as in
$\lesssim_\lambda$ or $\asymp_d$, indicate the parameters on which the
implicit constants may depend. Constants denoted by
$c,C,c_0,C_0,\dots$ may change from line to line, subject to any displayed
dependencies.

A well-specified finite mixture means that the data density is $p_{G_0}$
for a fixed, finitely supported mixing distribution $G_0$. We call a measure
$D$-atomic if it is supported on at most $D$ points. For a measure $G$,
$\supp(G)$ denotes its support; for a set $A$, $\conv(A)$ denotes its convex
hull; and $T_\#G$ denotes the pushforward of $G$ by a measurable map $T$.
We write $B(x,r)$ for a Euclidean ball, $x^+=\max\{x,0\}$,
$\log_+x=\max\{\log x,0\}$, and $\mathbf 1_A$ for the indicator of $A$.
The operator norm is $\|\cdot\|_{\op}$. An $L^2$ norm is written
$\|\cdot\|_2$ when its reference measure is clear and $\|\cdot\|_{L^2(Q)}$
when it is not. For functions on $\mathbb N_0=\{0,1,2,\ldots\}$,
$\ell^q$ and $\ell^q(Q)$ denote the norms with respect to counting measure
and a probability measure $Q$, respectively. If a maximizer has infinite
support, our support lower bounds hold automatically; the proofs therefore
need only rule out maximizers with too few atoms.

We also fix the bracketing notation used in the likelihood bounds. If
\(\mathcal F\) is a class of square-integrable functions with reference
measure \(Q\), an \(\eta\)-bracket is a pair \(\ell\le f\le u\), with
inequalities understood \(Q\)-a.e., such that
\(\|u-\ell\|_{L^2(Q)}\le\eta\). The bracketing number
\(N_{[]}(\eta,\mathcal F,L^2(Q))\) is the minimum number of such brackets
needed to cover \(\mathcal F\). Near a true density \(p_0\), we use the
coordinates \(\sqrt{p/p_0}-1\). Their \(L^2(p_0)\) norm is
\(\|\sqrt p-\sqrt{p_0}\|_2\), the Hellinger distance under our convention.

\subsection*{Gaussian location mixtures}

Let $\phi_d(x)=(2\pi)^{-d/2}e^{-|x|^2/2}$, $x\in\R^d$. For a probability measure $G$ on $\R^d$, set $p_G(x)=\int \phi_d(x-a)\,dG(a)$. Given $X_1,\dots,X_n\in\R^d$, an NPMLE is any maximizer of $G\mapsto \sum_{i=1}^n \log p_G(X_i)$ over probability measures on $\R^d$.

Our first theorem answers the support question for every fixed finite Gaussian
location mixture. The constant in the lower bound may depend on the true
mixture and on the dimension, but the power of \(\log n\) depends only on the
dimension.

\begin{theorem}
\label{thm:gaussian-finite-main}
Fix $d\ge1$, and let $G_0$ be a probability measure on $\R^d$ with finite nonempty support. Let $X_1,\dots,X_n$ be i.i.d.\ with density $p_{G_0}$, and let $\wh G_n$ be any Gaussian location NPMLE over $\R^d$. Then there exists a constant $c_{G_0,d}>0$ such that
\[
    \mathbb P\Bigl\{
        |\supp(\wh G_n)|
        \ge
        c_{G_0,d}\,
        \frac{(\log n)^{d/2}}{\log\log n}
    \Bigr\}\to1.
\]
\end{theorem}

For \(d=1\) and \(G_0=\delta_0\), the theorem gives
\(c\sqrt{\log n}/\log\log n\) fitted atoms from data drawn from a single
standard Gaussian. This gives the first diverging lower bound in the setting
asked about by Polyanskiy and Wu.

The proof compares the best likelihood gains available with and without a
support constraint. To lower-bound the unrestricted likelihood, first suppose
that \(G_0=\delta_0\). Moving small amounts of mass from zero to many separated
locations produces nearly uncorrelated likelihood scores and a total gain of
order \((\log n)^{d/2}\). For a general finite mixture, we perform the same
perturbation at a vertex of the convex hull of its support, moving mass outward
in directions for which that component dominates the other components. To
upper-bound the likelihood over all \(D\)-atomic mixtures, a localized
bracketing argument gives a gain of order at most \(D\log\log n\). Comparing
the two bounds yields the theorem.

\subsection*{Poisson mixtures}

Poisson mixtures also acquire extra atoms under a fixed finite mixture, but
the lower-bound rate is different. For \(\theta\ge0\), let
\(p_\theta(k)=e^{-\theta}\theta^k/k!\), \(k\in\mathbb N_0\), with
\(p_0\) interpreted as the point mass at zero. For a probability measure
\(G\) on \([0,\infty)\), write
\(p_G(k)=\int p_\theta(k)\,dG(\theta)\); we use the same notation for finite
nonnegative measures. Given observations in \(\mathbb N_0\), a
Poisson-mixture NPMLE maximizes \(\sum_i\log p_G(X_i)\) over all such
probability measures \(G\).

The Poisson theorem uses the largest true mean as the exposed component. The assumption that this largest mean is positive excludes only the degenerate law concentrated at zero.

\begin{theorem}
\label{thm:poisson-finite-main}
Let $G_0$ be a probability measure on $[0,\infty)$ with finite support, and suppose its largest support point $\lambda_*$ is positive. Let $X_1,\dots,X_n$ be i.i.d.\ with mass function $p_{G_0}$, and let $\wh G_n^{\mathrm{Pois}}$ be any Poisson-mixture NPMLE. Then there exists $c_{G_0}>0$ such that
\[
    \mathbb P\Bigl\{
        |\supp(\wh G_n^{\mathrm{Pois}})|
        \ge
        c_{G_0}\,\frac{\sqrt{\log n}}{(\log\log n)^{3/2}}
    \Bigr\}\to1.
\]
\end{theorem}

To construct a likelihood gain, we move mass from the largest true mean
\(\lambda_*\) to points
\(\lambda_*+a\sqrt{\lambda_*}\). After normalization, the resulting scores
have the same covariance kernel as the Gaussian scores. Their untruncated
fourth moments are dominated by count values far larger than any observation
in the sample and would restrict the construction to
\(|a|\lesssim(\log n)^{1/4}\). Truncating those irrelevant tails allows
\(a^2\le\gamma\log n/\log\log n\) for every fixed \(\gamma<1\), which
produces the stated rate.

In both one-dimensional models, known upper bounds are logarithmic. A fixed
finite Gaussian mixture is subgaussian, so Polyanskiy and Wu's
self-regularization theorem gives
\(|\supp(\wh G_n)|=O_{\mathbb P}(\log n)\)
\cite{polyanskiy2020self}. For Poisson mixtures, let
\(M_n=\max_iX_i\). Since every observed count lies in
\(\{0,\dots,M_n\}\), Lindsay's theorem gives a maximizer with at most
\(M_n+1\) atoms \cite{lindsay1983geometryA}; uniqueness then gives the same
bound for the NPMLE itself \cite{lindsay1993uniqueness}. Under every fixed
finite Poisson mixture with positive largest mean,
\(M_n=(1+o_{\mathbb P}(1))\log n/\log\log n\). In both models, therefore,
the lower and upper bounds differ by a factor of order \(\sqrt{\log n}\),
apart from powers of \(\log\log n\).

For Gaussian mixtures in dimensions \(d\ge2\), the gap is much larger. The
only general comparison is the finite-sample bound that some maximizer has at
most \(n\) atoms. No upper bound is known near the polylogarithmic lower bound
\((\log n)^{d/2}/\log\log n\), and nonuniqueness prevents the bound for one
selected maximizer from applying automatically to all maximizers.

\subsection*{Related literature and empirical Bayes motivation}

In empirical Bayes applications, the NPMLE mixing distribution is used as an
estimated prior; see, for example, Jiang--Zhang \cite{jiang2009general},
Koenker--Mizera \cite{koenker2014convex}, Dicker--Zhao
\cite{dicker2016high}, and Soloff--Guntuboyina--Sen
\cite{soloff2024multivariate}. This use raises two separate questions: how
well the fitted mixture estimates the data distribution or Bayes rule, and how
many atoms the maximizing mixing distribution contains. Our results address
the second question under a fixed finite-mixture data law.

Previous work on NPMLE support size has emphasized computation and upper
bounds. Koenker and Gu's \emph{Minimalist $g$-Modeling}
\cite{koenker2019minimalist} documented sparse NPMLE fits; as reported by
Polyanskiy and Wu \cite{polyanskiy2020self}, their numerical examples suggested
support sizes on the \(\sqrt n\) scale. Polyanskiy and Wu proved the much
smaller upper bound \(|\supp(\wh G_n)|=O_{\mathbb P}(\log n)\) under
subgaussian sampling and asked whether it is tight for data from a single
Gaussian component. Polyanskiy and Sellke \cite{polyanskiySellke2025} study
certified computation and support behavior for Gaussian-location NPMLEs.

A separate literature establishes statistical guarantees for mixture NPMLEs.
Zhang \cite{zhang2009generalized} and Saha--Guntuboyina
\cite{saha2020nonparametric} prove density-risk bounds for Gaussian mixture
NPMLEs, while Doss--Wu--Yang--Zhou \cite{doss2020optimal} treat
high-dimensional Gaussian location mixtures. Zhang--Volgushev
\cite{zhangVolgushev2026adaptivity} prove adaptation to finitely supported
mixing distributions in Gaussian and Poisson models.
Ghosh--Guntuboyina--Mukherjee--Tran
\cite{GhoshGuntuboyinaMukherjeeTran2026} obtain KL stability bounds for
Gaussian-mixture NPMLEs and their approximate solutions, and Chen--Wu
\cite{ChenWu2026RegretHellinger} prove regret--Hellinger inequalities for
Gaussian empirical Bayes. These guarantees do not control the number of atoms
in the optimizer. The present results show that accurate estimation can
coexist with diverging fitted support under a fixed finite mixture.

Mixture likelihood-ratio statistics also have nonstandard asymptotics.
Hartigan \cite{hartigan1985failure} gave an early example for normal mixtures;
Aza{\"i}s--Gassiat--Mercadier \cite{azais2009likelihood} and Jiang--Zhang
\cite{jiangZhang2019rate} obtained further nonstandard limits; and
Zhang--Volgushev \cite{zhangVolgushev2025surprising} proved a
tight-versus-divergent dichotomy for nonparametric mixture likelihood-ratio
tests with bounded parameter spaces. Those works study the value of a
likelihood-ratio statistic. We instead study the smallest support size among
likelihood maximizers.

\section{Likelihood Geometry and Entropy Tools}
\label{sec:tools}

The proofs of Theorems~\ref{thm:gaussian-finite-main} and
\ref{thm:poisson-finite-main} use a common upper bound for mixtures with a
prescribed number of atoms. After constructing an unrestricted likelihood
gain, we must show that
no mixture with only \(D\) atoms can match it. This section supplies the two
reductions used for that comparison. The first restricts the atoms of a
competitor to a sample-dependent compact set without decreasing likelihood;
the second uses
localized bracketing entropy to bound the likelihood gain over the resulting
sparse class. The model-specific applications appear in
Propositions~\ref{prop:gaussian-finite-upper} and
\ref{prop:poisson-finite-upper}.

\subsection{Compact Support Reduction}
\label{sec:compact-support-reduction}

For Gaussian mixtures, atoms outside the convex hull of the observations can
be projected onto that hull without decreasing the fitted density at any
observation. This reduction depends only on the Gaussian kernel, not on the
distribution that generated the data.

\begin{lemma}
\label{lem:convex-support}
Let $C_n=\conv\{X_1,\dots,X_n\}\subset\R^d$. For every probability measure $G$ on $\R^d$, there exists a probability measure $\widetilde G$ supported on $C_n$ such that $p_{\widetilde G}(X_i)\ge p_G(X_i)$ for every $1\le i\le n$. Moreover, if $G$ has at most $D$ atoms, then so does $\widetilde G$.
\end{lemma}

\begin{proof}
Let $\Pi_{C_n}$ be Euclidean projection onto $C_n$, and let $\widetilde G=(\Pi_{C_n})_\#G$ be the pushforward of $G$. If $a\notin C_n$ and $y=\Pi_{C_n}(a)$, then $(a-y)\cdot(x-y)\le0$ for every $x\in C_n$.
Therefore
\begin{align*}
    \log \phi_d(x-y)-\log \phi_d(x-a)
    &=
    (y-a)\cdot x-\frac12(|y|^2-|a|^2) \\
    &=
    \frac12|a-y|^2+(a-y)\cdot(y-x)\ge0.
\end{align*}
Applying this with $x=X_i$ shows that replacing an atom at $a$ by an atom of the same mass at $y$ does not decrease any fitted density value $p_G(X_i)$. Integrating this pointwise inequality over $a$ gives $p_{\widetilde G}(X_i)\ge p_G(X_i)$ for every $i$.
\end{proof}

\subsection{Localized Bracketing Bound}
\label{sec:localized-bracketing}

For every candidate density \(p\), its expected log-likelihood relative to
the truth is the negative Kullback--Leibler divergence
\(-\mathrm{KL}(p_0,p)\). Candidates far from \(p_0\) therefore have a
larger negative mean and need not be covered as
finely as candidates near \(p_0\). Local bracketing formalizes this idea by
bounding the entropy separately on each Hellinger ball. We use the following
localized form of Wong--Shen
\cite[Theorem~1, the remark following it, and Lemma~7]{wong1995probability};
see also van de Geer \cite[Chapter~7]{vandegeer2000empirical}. The theorem
controls the likelihood outside a Hellinger ball, while the truncated-log
estimate in their Lemma~7 controls the innermost ball.

\begin{lemma}
\label{lem:wong-shen-local}
Let $X_1,\dots,X_n$ be i.i.d.\ with density $p_0$ with respect to a common dominating measure \(\nu\), and set \(P_0(dx)=p_0(x)\nu(dx)\). Let $\mathcal P_n$ be a class of densities \(p\) with respect to \(\nu\) such that \(p\,d\nu\ll P_0\); ratios \(p/p_0\) are understood \(P_0\)-a.e. Assume the suprema below are measurable; without this assumption, interpret all probabilities as outer probabilities. For $0<\delta\le\sqrt2$, define
\[
    \mathcal F_n(\delta)
    =
    \Bigl\{
        \sqrt{p/p_0}-1:
        p\in\mathcal P_n,\
        \|\sqrt{p/p_0}-1\|_{L^2(p_0)}\le\delta
    \Bigr\}.
\]
Assume \(A_n\to\infty\) and \(A_n=o(n)\). Fix \(C<\infty\), choose
\(K_{\mathrm{WS}}=K_{\mathrm{WS}}(C)\) sufficiently large, and set
\(\delta_n=(K_{\mathrm{WS}}A_n/n)^{1/2}\). If
\[
    J_{[],n}(\delta)
    :=
    \int_0^\delta
    \sqrt{
        1+\log N_{[]}\bigl(
            \eta,\mathcal F_n(\delta),L^2(p_0)
        \bigr)
    }
    \,d\eta
    \le C\delta\sqrt{A_n},
\]
for every $\delta_n\le\delta\le\sqrt2$, then there is a constant $C_1<\infty$, depending only on $C$, such that
\[
    \mathbb P\left\{
    \sup_{p\in\mathcal P_n}
    \sum_{i=1}^n\log\frac{p(X_i)}{p_0(X_i)}
    \le C_1A_n
    \right\}\to1.
\]
\end{lemma}

\begin{proof}
This is the standard localization of Wong--Shen's inequalities. Peel the
region outside the ball of radius \(\delta_n\) into dyadic Hellinger shells.
On the shell with outer radius \(\delta\), apply
\cite[Theorem~1]{wong1995probability} to the Hellinger ball of radius
\(2\delta\). The assumed entropy integral is at most
\(2C\delta\sqrt{A_n}\). Once \(K_{\mathrm{WS}}(C)\) is large enough, the
theorem's critical-radius condition holds, and the probability that any
density in this shell has log-likelihood ratio larger than
\(-c n\delta^2\) is at most \(C'e^{-c'n\delta^2}\). Summing over the shells
costs \(O(e^{-cA_n})\).

It remains to control the ball of radius \(\delta_n\). Apply the
truncated-log estimate in \cite[Lemma~7]{wong1995probability}, with the
truncation level chosen as in the remark following their Theorem~1. The same
entropy integral bounds the centered truncated empirical process by
\(C_1A_n\), while the truncation remainder has probability
\(O(e^{-cA_n})\). Since \(A_n\to\infty\) and \(A_n=o(n)\), both applications
are in their stated range. Combining the inner ball with the dyadic shells
proves the lemma.
\end{proof}

In the Gaussian application, \(A_n\) is proportional to the atomic budget
times \((d+1)\log\log n\); in the Poisson application, it is proportional to
the atomic budget times \(\log\log n\). In both cases, the factor
\(\log\log n\) comes from the logarithm of the sample-dependent parameter
radius.

Both entropy proofs decompose a candidate mixture into a baseline submixture
near the true atoms and several additional components. The next lemma controls
how their square-root density increments combine. Recording this algebra once
leaves only model-specific tail estimates in the Gaussian and Poisson
sections.

\begin{lemma}
\label{lem:contaminant-calculus}
On a measure space with dominating measure \(\nu\), let \(p_0\) be a
positive density, let \(a\) be a positive subdensity, and let \(q\) be a
density, all with respect to \(\nu\). For \(w>0\), set
\[
 h_{a,w,q}=\frac{\sqrt{a+wq}-\sqrt a}{\sqrt{p_0}},
 \qquad e_{a,w,q}=\|h_{a,w,q}\|_{L^2(p_0)},
 \qquad J_{a,w,q}=\int\frac{w^2q^2}{a+wq}\,d\nu.
\]
In this lemma, \(L^2(p_0)\) abbreviates \(L^2(p_0\,d\nu)\).
Then \(J_{a,w,q}/4\le e_{a,w,q}^2\le J_{a,w,q}\), and, pointwise,
\[
 \frac12h_{a,w,q}
 \le\partial_{\log w}h_{a,w,q}\le h_{a,w,q},
 \qquad
 \frac12\le\partial_{\log w}\log e_{a,w,q}\le1.
\]
If \(s=\sqrt{a/p_0}\), then replacing \(s\) by any nonnegative measurable
function \(s_0\) changes
\(h_{a,w,q}\) by at most \(|s-s_0|\) pointwise.

More generally, if
\(p=a+\sum_{j=1}^Jw_jq_j\) and \(h_j=h_{a,w_j,q_j}\), then
\[
 \sqrt{p/p_0}=\Phi(s,h_1,\ldots,h_J),
 \qquad
 \Phi(s,h_1,\ldots,h_J)
 =\left\{s^2+\sum_{j=1}^J(h_j^2+2sh_j)\right\}^{1/2}.
\]
On the nonnegative orthant, \(\Phi\) is coordinatewise increasing and
\(\partial_{h_j}\Phi\le1\), \(\partial_s\Phi\le\sqrt{J+1}\).
\end{lemma}

\begin{proof}
The identity
\[
 \sqrt{a+wq}-\sqrt a
 =\frac{wq}{\sqrt{a+wq}+\sqrt a}
\]
gives the comparison with \(J_{a,w,q}\). Differentiating in \(\log w\)
shows that the ratio of \(\partial_{\log w}h_{a,w,q}\) to
\(h_{a,w,q}\) is
\((\sqrt{a+wq}+\sqrt a)/(2\sqrt{a+wq})\in[1/2,1]\); integrating against
\(h_{a,w,q}^2p_0\) gives the bound for \(\partial_{\log w}\log e_{a,w,q}\).
The map \(x\mapsto\sqrt{x^2+c}-x\) is one-Lipschitz on
\([0,\infty)\), which proves the baseline comparison.

Finally, \(w_jq_j/p_0=h_j^2+2sh_j\), proving the composition identity.
Its derivatives are
\[
 \partial_{h_j}\Phi=(s+h_j)/\Phi\le1,
 \qquad
 \partial_s\Phi=(s+\sum_jh_j)/\Phi\le\sqrt{J+1},
\]
the last inequality following from Cauchy--Schwarz.
\end{proof}


\section{Gaussian Support Divergence}
\label{sec:gaussian-support}

This section proves Theorem~\ref{thm:gaussian-finite-main}. The proof compares
the two likelihood levels established in
Propositions~\ref{prop:gaussian-finite-constructive} and
\ref{prop:gaussian-finite-upper}: the first constructs an unrestricted gain of
order \((\log n)^{d/2}\), while the second bounds the gain of every
\(D\)-atomic mixture by \(O(D\log\log n)\). Their comparison forces the
support of the NPMLE to diverge at the stated rate.

Fix a finite mixing distribution
\[
    G_0=\sum_{m=1}^K\pi_m\delta_{\mu_m},
    \qquad \pi_m>0,
\]
on \(\R^d\), after merging duplicate atoms.  Let \(p_0=p_{G_0}\), and write likelihoods relative to \(p_0\):
\[
    L_n(G;G_0)=\sum_{i=1}^n\log\frac{p_G(X_i)}{p_0(X_i)}.
\]
The perturbation moves a small amount of mass from one true atom \(\mu_*\) to
many points \(\mu_*+a\). We choose \(\mu_*\) to be a vertex of the convex
hull of the true support and restrict \(a\) to a cone pointing away from the
other atoms. In the Gaussian integrals that determine the score moments, this
geometry makes the \(\mu_*\)-component dominate the mixture density. The
scores are then comparable to those of a single Gaussian component and become
nearly uncorrelated when their displacement vectors are separated.

\subsection{Scores From an Exposed Component}
\label{sec:gaussian-exposed-scores}

The constructive half of the proof requires many nearly orthogonal score
directions. We obtain them by perturbing an exposed atom of the true mixing
distribution outward from the convex hull of the remaining atoms; the lemmas
in this subsection provide the normalization, correlation decay, and moment
bounds used in Proposition~\ref{prop:gaussian-finite-constructive}.

Suppose first that \(K>1\). Choose an index \(m_*\) for which
\(\mu_{m_*}\) is a vertex of
\(\conv\{\mu_1,\dots,\mu_K\}\), and write
\(\mu_*=\mu_{m_*}\) and \(\pi_*=\pi_{m_*}\). A sufficiently narrow cone
inside the outward normal cone at \(\mu_*\) has nonempty interior and satisfies
\begin{equation}
\label{eq:exposed-cone}
    (\mu_*-\mu_m)\cdot a\ge \eta_0 |a|,
    \qquad a\in\mathcal C,
    \quad m:\mu_m\ne\mu_* ,
\end{equation}
for some \(\eta_0>0\). We take the cone narrow enough that, for all
\(a,b\in\mathcal C\),
\[
    |a+b|\ge c_{\mathcal C}(|a|+|b|),
    \qquad a\cdot b\ge0.
\]
If \(K=1\), there are no competing true atoms. We set \(m_*=1\),
\(\mu_*=\mu_1\), \(\pi_*=1\), and choose any closed cone with nonempty
interior on which \(|a+b|\ge c_{\mathcal C}(|a|+|b|)\) and
\(a\cdot b\ge0\).

For \(a\in\mathcal C\), define
\[
    \Delta_a(x)=\phi_d(x-\mu_*-a)-\phi_d(x-\mu_*),
    \qquad
    g_a(x)=\frac{\Delta_a(x)}{p_0(x)},
\]
\[
    \sigma_a^2=\E_{p_0}g_a(X)^2
    =\int \frac{\Delta_a(x)^2}{p_0(x)}\,dx,
    \qquad
    h_a=\frac{g_a}{\sigma_a}.
\]
Then \(\E_{p_0}h_a(X)=0\) and \(\E_{p_0}h_a(X)^2=1\).

The next lemma supplies the three estimates needed for the likelihood
construction. The first identifies the normalization of each score, the
second shows that separated displacement vectors have weakly correlated
scores, and the third provides the moments needed for concentration. All
three estimates match the corresponding one-component formulas up to
constants depending on \(G_0\).

\begin{lemma}
\label{lem:gaussian-exposed-scores}
There are constants \(R_0<\infty\) and \(0<c<C<\infty\), depending only on \(G_0,d\) and \(\mathcal C\), such that, uniformly for \(a,b\in\mathcal C\) with \(|a|,|b|\ge R_0\),
\[
    c e^{|a|^2}\le \sigma_a^2\le C e^{|a|^2},
\]
\[
    \bigl|\E_{p_0}h_a(X)h_b(X)\bigr|
    \le
    C e^{-|a-b|^2/2},
    \qquad a\ne b,
\]
and, for each fixed integer \(m\ge3\),
\[
    \E_{p_0}|h_a(X)|^m
    \le
    C_m e^{m(m-2)|a|^2/2}.
\]
\end{lemma}

\begin{proof}
Translate so that \(\mu_*=0\). Then
\[
    p_0(y)=\phi_d(y)
    \left[
        \pi_*+
        \sum_{m\ne m_*}\pi_m
        \exp\{\mu_m\cdot y-|\mu_m|^2/2\}
    \right]
    =\phi_d(y)W(y).
\]
By \eqref{eq:exposed-cone}, \(\mu_m\cdot a\le-\eta_0|a|\) for \(a\in\mathcal C\) and \(m\ne m_*\). Under the tilted density proportional to \(\phi_d(y-a)^2/\phi_d(y)\), the variable \(Y\) has law \(\cN(2a,\Id_d)\) and
\[
    \frac{\phi_d(Y-a)^2}{\phi_d(Y)}=e^{|a|^2}\phi_d(Y-2a).
\]
On the event \(|Y-2a|\le c|a|\), every competing component in \(W(Y)\) is \(O(e^{-c|a|})\), while the complement has probability \(e^{-c|a|^2}\). Hence
\[
    \int \frac{\phi_d(y-a)^2}{p_0(y)}\,dy
    =
    \pi_*^{-1}e^{|a|^2}\{1+O(e^{-c|a|})\}.
\]
The unshifted-square and cross terms in \(\Delta_a^2\) are both \(O(1)\),
whereas the displayed integral has order \(e^{|a|^2}\). This proves
\(\sigma_a^2\asymp e^{|a|^2}\).

For covariance, the leading integral is
\[
    \int\frac{\phi_d(y-a)\phi_d(y-b)}{p_0(y)}\,dy
    =
    e^{a\cdot b}\,\E\frac1{W(Y_{a,b})},
    \qquad
    Y_{a,b}\sim\cN(a+b,\Id_d).
\]
Since \(a+b\in\mathcal C\) and
\(|a+b|\ge c_{\mathcal C}(|a|+|b|)\), the exposed-component domination
used for the variance also gives
\[
    \E\frac1{W(Y_{a,b})}
    =
    \pi_*^{-1}\{1+O(e^{-c(|a|+|b|)})\}.
\]
Dividing by \(\sigma_a\sigma_b\asymp e^{(|a|^2+|b|^2)/2}\) gives the displayed \(e^{-|a-b|^2/2}\) term, with a smaller domination error. The remaining three terms in \(\Delta_a\Delta_b\) contain at least one unshifted factor \(\phi_d(y)\). After normalization they are bounded by \(C e^{-(|a|^2+|b|^2)/2}\), which is at most \(C e^{-|a-b|^2/2}\) because \(a\cdot b\ge0\) on the cone.

For moments, use \(p_0\ge\pi_*\phi_d(\cdot-\mu_*)\) and \(|\Delta_a|^m\le2^{m-1}\{\phi_d(\cdot-\mu_*-a)^m+\phi_d(\cdot-\mu_*)^m\}\). The Gaussian identity
\[
    \int \phi_d(y-a)^m\phi_d(y)^{1-m}\,dy
    =
    \exp\{m(m-1)|a|^2/2\}
\]
combined with \(\sigma_a^m\ge c_m e^{m|a|^2/2}\) proves the moment bound.
\end{proof}

\subsection{Constructive Likelihood Gain}
\label{sec:gaussian-constructive-gain}

We now use the exposed scores to prove the constructive half of
Theorem~\ref{thm:gaussian-finite-main}. A separated grid of perturbations
supplies order \((\log n)^{d/2}\) nearly orthogonal directions, and optimizing
their small mixing weights yields the gain stated in
Proposition~\ref{prop:gaussian-finite-constructive}.

Let \(L=\log n\), and fix a small constant \(A_*=A_*(G_0,d)>0\). Set
\[
    K_{n,d}^{G_0}
    =
    \{a\in\mathcal C:R_0\le |a|\le A_*\sqrt L\}.
\]
Choose \(L_d\) large, and let \(\mathcal A_{n,d}\) be a maximal \(L_d\)-separated subset of \(K_{n,d}^{G_0}\). Since \(\mathcal C\) has positive solid angle,
\[
    J_{n,d}:=|\mathcal A_{n,d}|\asymp_{G_0,d}(\log n)^{d/2}.
\]

The following proposition gives the unrestricted half of the likelihood
comparison. Its sample-dependent mixing distribution removes a small amount
of mass from \(\mu_*\) and redistributes it over
\(\mu_*+\mathcal A_{n,d}\). We use this distribution only as a feasible
competitor; it is not asserted to be the NPMLE.

\begin{proposition}
\label{prop:gaussian-finite-constructive}
For suitable constants \(R_0,L_d,A_*\), there is a random mixing distribution \(G_n\) such that
\[
    \mathbb P\left\{
        L_n(G_n;G_0)
        \ge c_{G_0,d}J_{n,d}
    \right\}\to1.
\]
\end{proposition}

\begin{proof}
For \(a\in\mathcal A_{n,d}\), let
\[
    d_a=\frac1{\sqrt n}\sum_{i=1}^n h_a(X_i),
    \qquad
    \wh H_{ab}=\frac1n\sum_{i=1}^n h_a(X_i)h_b(X_i),
    \qquad
    H_{ab}=\E_{p_0}h_a(X)h_b(X).
\]
Thus \(d_a\) is the empirical score in direction \(a\), while \(\wh H\)
and \(H\) are its empirical and population Gram matrices.
By Lemma~\ref{lem:gaussian-exposed-scores}, choosing \(L_d\) large makes \(\sup_a\sum_{b\ne a}|H_{ab}|<1/5\). Gershgorin's theorem gives \(\lambda_{\max}(H)\le1.2\). The moment estimate in Lemma~\ref{lem:gaussian-exposed-scores}, with \(A_*\) small enough, gives
\[
    \E |h_a(X)|^3\le C n^{3A_*^2/2},
    \qquad
    \E |h_a(X)|^4\le C n^{4A_*^2}.
\]
Since \(J_{n,d}\) is polylogarithmic, Chebyshev's inequality and a union
bound over the \(J_{n,d}^2\) entries, followed by
\(\|A\|_{\op}\le J_{n,d}\max_{a,b}|A_{ab}|\), give
\[
    \|\wh H-H\|_{\op}=o_{\mathbb P}(1).
\]
We fix \(A_*\) small enough that the resulting powers of \(n\) are
negligible in the concentration estimates that follow; any sufficiently small
constant satisfying \(A_*^2<1/8\) suffices.

Let \(S=\sum_{a\in\mathcal A_{n,d}}(d_a^+)^2\). We next show that a positive
fraction of the score coordinates contribute. Fix a large constant \(B\) and
write \(\psi_B(x)=(x^+)^2\wedge B\). The one- and two-dimensional
Berry--Esseen bounds for bounded Lipschitz test functions, together with the
third-moment estimate in Lemma~\ref{lem:gaussian-exposed-scores}, hold
uniformly over all coordinates and pairs;
their error is \(n^{-c}\) for some \(c>0\). Consequently,
\[
    \E\psi_B(d_a)=\E\psi_B(Z)+o(1)
\]
uniformly in \(a\), where \(Z\sim\mathcal N(0,1)\). For a centered bivariate
normal pair \((Z_a,Z_b)\) with correlation \(H_{ab}\), the Gaussian maximal
correlation inequality gives
\[
    |\Cov(\psi_B(Z_a),\psi_B(Z_b))|
    \le C_B|H_{ab}|.
\]
The uniform bivariate approximation and the row-sum bound for \(H\) therefore
imply
\[
    \Var\left(\sum_a\psi_B(d_a)\right)
    =O_B(J_{n,d})+o(J_{n,d}^2).
\]
Choose \(B\) so that \(\E\psi_B(Z)>1/3\). Chebyshev's inequality and
\(\psi_B(d_a)\le(d_a^+)^2\) then show that \(S\) is at least a constant
multiple of \(J_{n,d}\) with high probability. The identity
\(\E\sum_a d_a^2=J_{n,d}\) gives the matching stochastic upper bound:
\[
    \mathbb P\{S\ge cJ_{n,d}\}\to1,
    \qquad
    S=O_{\mathbb P}(J_{n,d}).
\]

Fix \(\eta=1/5\), put \(r_a=d_a^+/(1+\eta)\), and set
\[
    t_a=\frac{r_a}{\sqrt n\,\sigma_a}.
\]
Since \(\inf_a\sigma_a>0\) and \(\sum_a r_a\le\sqrt{J_{n,d}S}=O_{\mathbb P}(J_{n,d})=o_{\mathbb P}(\sqrt n)\), we have \(\sum_a t_a=o_{\mathbb P}(1)\). Hence \(\sum_a t_a\le\pi_*\) with probability tending to one. On this event define
\[
    G_n=
    G_0-
    \Bigl(\sum_{a\in\mathcal A_{n,d}}t_a\Bigr)\delta_{\mu_*}
    +
    \sum_{a\in\mathcal A_{n,d}}t_a\delta_{\mu_*+a},
\]
and set \(G_n=G_0\) otherwise.

On the event \(\max_i|X_i|\le C_{G_0}\sqrt L\), which has probability tending to one, the bound \(p_0\ge\pi_*\phi_d(\cdot-\mu_*)\) gives
\[
    |h_a(X_i)|
    \le
    C\exp\{a\cdot(X_i-\mu_*)-|a|^2\}+C
    \le n^\rho
\]
for all \(i\) and \(a\in\mathcal A_{n,d}\), where \(\rho<1/2\) after choosing \(A_*\) small enough. Therefore, with
\[
    u_i=\frac1{\sqrt n}\sum_{a\in\mathcal A_{n,d}}r_a h_a(X_i),
\]
we have \(\max_i|u_i|=o_{\mathbb P}(1)\), and
\[
    \frac{p_{G_n}(X_i)}{p_0(X_i)}=1+u_i.
\]
Moreover, \(\sum_i u_i^2=r^\top\wh H r=O_{\mathbb P}(J_{n,d})\).
The remainder in \(\log(1+u)=u-u^2/2+O(|u|^3)\) is therefore at most
\(O(\max_i|u_i|\sum_i u_i^2)=o_{\mathbb P}(J_{n,d})\), and hence
\[
    L_n(G_n;G_0)
    =
    r^\top d-\frac12 r^\top\wh H r+o_{\mathbb P}(J_{n,d}).
\]
Since \(r^\top d=S/(1+\eta)\) and
\(r^\top\wh H r\le(1.2+o_{\mathbb P}(1))S/(1+\eta)^2\), the resulting
coefficient of \(S\) is positive for \(\eta=1/5\). The high-probability
lower bound on \(S\) proves the proposition.
\end{proof}

\subsection{Localized Upper Bound}
\label{sec:gaussian-localized-upper}

We now bound the likelihood of every \(D\)-atomic competitor. Lemma
~\ref{lem:convex-support} restricts its atoms to a ball whose radius grows
like \(\sqrt{\log n}\). To retain the desired \(D\log\log n\) likelihood
scale, the entropy must therefore depend on this radius only through its
logarithm.

For \(D\ge1\), \(M\ge2\), and \(0<\delta\le\sqrt2\), let \(\mathcal P_{d,D}(M)\) be the class of Gaussian location mixtures with at most \(D\) atoms supported in \(B(0,M)\), and define
\[
    \mathcal F_{d,D}^{G_0}(M,\delta)
    =
    \left\{
        \sqrt{p_G/p_0}-1:
        p_G\in\mathcal P_{d,D}(M),
        \ \|\sqrt{p_G/p_0}-1\|_{L^2(p_0)}\le\delta
    \right\}.
\]

The next lemma provides this dependence. The factor \(\log M\), rather than
a power of \(M\), is what becomes \(\log\log n\) when
\(M\asymp\sqrt{\log n}\).

\begin{lemma}
\label{lem:gaussian-finite-local-bracket}
For fixed \(G_0\), fixed \(d\), and fixed \(C_0<\infty\), there exists
\(C_{G_0,d,C_0}<\infty\) such that, for every \(M\ge2\),
\(1\le D\le M^{C_0}\), and \(0<\eta\le\delta\le\sqrt2\),
\[
    \log N_{[]}
    \bigl(
        \eta,
        \mathcal F_{d,D}^{G_0}(M,\delta),
        L^2(p_0)
    \bigr)
    \le
    C_{G_0,d,C_0}D(d+1)\bigl(\log M+\log(\delta/\eta)\bigr).
\]
\end{lemma}

\begin{proof}
Write \(\rho_G=\|\sqrt{p_G/p_0}-1\|_2\). We first work in a fixed
neighborhood \(0<\rho_G\le\rho_0\). Choose disjoint bounded neighborhoods
\(A_1,\ldots,A_K\) of the atoms of \(G_0\), and let \(A_0\) be their
complement. The local-geometry theorem of Gassiat--van Handel
\cite[Theorem~3.10]{gassiatVanHandel2014} controls both the mass outside
these neighborhoods and the perturbation formed by the remaining atoms,
uniformly in the number of candidate atoms. In particular,
\begin{equation}
\label{eq:gaussian-local-geometry}
 W:=G(A_0)\le C\rho_G,
 \qquad
 \left\|\sqrt{p_{G|_{A_0^c}}/p_0}-1\right\|_2\le C\rho_G.
\end{equation}
The theorem directly bounds \(W\). Its coefficient lower bound, followed by
second-order Taylor expansion of the kernels on the fixed sets \(A_i\), also
gives the second inequality. The same coefficient and Taylor bounds give the
pointwise envelope needed below for the innermost Hellinger ball: the
normalized inner density perturbation is bounded by a fixed \(L^2(p_0)\)
function formed from the Gaussian kernel and its first two location
derivatives on \(\bigcup_iA_i\).

Fix the contribution \(a\) of the atoms in \(\bigcup_iA_i\). We first
construct brackets after adding one atom outside these neighborhoods; the
full mixture will follow by composing at most \(D\) such brackets.
Write \(q_\theta=\phi_d(\cdot-\theta)\), and use the notation
\(h_{a,w,q_\theta}\), \(e_{a,w,q_\theta}\) from Lemma
~\ref{lem:contaminant-calculus}. Differentiation gives
\[
 \nabla_\theta h_{a,w,q_\theta}
 =\frac{wq_\theta(x-\theta)}
 {2\sqrt{p_0}\sqrt{a+wq_\theta}}.
\]
Because \(a\) contains a fixed positive amount of Gaussian mass with
centers in a fixed ball, the measure governing the derivative norm has a
uniform exponential radial moment. More precisely, if
\(d\mu=w^2q_\theta^2(a+wq_\theta)^{-1}dx\), comparison along rays gives
\begin{equation}
\label{eq:gaussian-harmonic-moment}
 \int e^{c|x|/(1+M)}
 \left(1+\frac{|x|}{1+M}\right)^2d\mu
 \le C\int d\mu,
 \qquad |\theta|\le M.
\end{equation}
To prove \eqref{eq:gaussian-harmonic-moment}, put
\(\widetilde q=wq_\theta\),
\(v=\widetilde q^2/a\), and
\(\mu=\widetilde qv/(\widetilde q+v)\); this last expression equals the
density in \eqref{eq:gaussian-harmonic-moment}. If \(x=t\omega\), with
\(|\omega|=1\), then for \(r\ge0\),
\[
 \widetilde q((t+r)\omega)
 \le e^{-r(t-M)-r^2/2}\widetilde q(t\omega),
 \qquad
 a((t+r)\omega)
 \ge e^{-r(t+B_0)-r^2/2}a(t\omega),
\]
where \(B_0\) bounds the inner centers. Hence
\(v((t+r)\omega)\le
e^{-r(t-2M-B_0)-r^2/2}v(t\omega)\). The harmonic-mean map
\((x,y)\mapsto xy/(x+y)\) is increasing and homogeneous, so the same
bound holds for \(\mu\). After polar integration, if
\(F(t)=t^{d-1}\int_{\mathbb S^{d-1}}\mu(t\omega)\,d\omega\), then
\[
 F(t+1)\le e^{-t/3}F(t),
 \qquad t\ge C(1+M+d).
\]
Iterating this inequality on unit shells proves
\eqref{eq:gaussian-harmonic-moment}. Lemma
~\ref{lem:contaminant-calculus} now yields
\[
 \|\nabla_\theta h_{a,w,q_\theta}\|_2
 \le C(1+M)e_{a,w,q_\theta}.
\]

We use these derivative bounds to construct ordered pointwise brackets,
which is stronger than constructing an \(L^2\)-cover. Partition the location
ball into boxes of diameter \(c/(1+M)\), and partition the positive weights
on the \(\log w\) scale. Within one location box, the ratio of any kernel to
the kernel at its center is at most
\(\exp\{c_1(1+|x|/(1+M))\}\), where \(c_1\) can be made smaller than the
constant in \eqref{eq:gaussian-harmonic-moment}. After squaring, that estimate
and \eqref{eq:gaussian-harmonic-moment} bound the \(L^2(p_0)\) norms of the
pointwise suprema of the weight and location derivatives by
\(Ce\) and \(C(1+M)e\), respectively.

Counting the weight and location cells gives, uniformly over the admissible
inner subdensities,
\begin{equation}
\label{eq:gaussian-one-contaminant-bracket}
 \log N_{[]}\left(
 \epsilon,
 \{h_{a,w,q_\theta}:|\theta|\le M,\ e_{a,w,q_\theta}\le r\},
 L^2(p_0)\right)
 \le Cd\{\log M+\log_+(r/\epsilon)\}.
\end{equation}
To see the endpoint \(w=0\), fix a preliminary location box and a point
\(\theta_0\) in it. By Lemma~\ref{lem:contaminant-calculus}, the logarithmic
derivative of \(e_{a,w,q_{\theta_0}}\) lies in \([1/2,1]\). All weights for
which this norm is at most \(c\epsilon\) therefore fit into one bracket
\([0,\sup_{\theta\text{ in the same box}}h_{a,w_*,q_\theta}]\) of width
\(C\epsilon\); above that
threshold, the relevant \(\log w\)-interval has length
\(O(1+\log(r/\epsilon))\). Meshes of size \(c\epsilon/r\) in \(\log w\)
and \(c\epsilon/[r(1+M)]\) in location complete the count. Adjusting the
mesh constants gives width \(\epsilon\).

We now bracket a Hellinger shell \(r/2<\rho_G\le r\). Decompose
\[
 p_G=a+\sum_{j=1}^Jw_jq_{\theta_j},\qquad J\le D,
\]
where \(a=p_{G|_{A_0^c}}\). By \eqref{eq:gaussian-local-geometry}, with
\(s=\sqrt{a/p_0}\),
\(\|s-1\|_2\le Cr\). Positivity also gives
\[
 \|h_{a,w_j,q_{\theta_j}}\|_2
 \le\|\sqrt{p_G}-\sqrt a\|_2\le Cr.
\]
Let \(\mathcal S_D^{\rm in}(r)\) denote the class of these inner functions:
\[
 \begin{split}
 \mathcal S_D^{\rm in}(r)=\bigl\{\sqrt{p_H/p_0}:{}&
 H\text{ is a subprobability measure},\quad |\supp(H)|\le D,\\
 &\supp(H)\subset\textstyle\bigcup_iA_i,\quad
 \|\sqrt{p_H/p_0}-1\|_2\le Cr\bigr\}.
 \end{split}
\]
For \(a=p_H\), write \(a=m\bar a\), where \(m=H(\R^d)\) and
\(\bar a\) is a probability-mixture density. The reverse triangle
inequality gives \(|\sqrt m-1|\le Cr\), and hence
\[
 \|\sqrt{\bar a/p_0}-1\|_2
 \le m^{-1/2}\|\sqrt{a/p_0}-1\|_2+|m^{-1/2}-1|
 \le Cr.
\]
Apply Gassiat--van Handel's fixed-compact local bracketing theorem to
\(\bar a\), with order \(D\vee K\), and bracket the scalar \(\sqrt m\)
separately. Multiplying the nonnegative scalar and density endpoints gives
\[
 \log N_{[]}\bigl(\epsilon,
 \mathcal S_D^{\rm in}(r),L^2(p_0)\bigr)
 \le CD(d+1)\log(CDr/\epsilon).
\]

Put \(\zeta=c\min\{\eta,r\}\). Bracket \(s\) at width
\(\zeta/(CD)\), choose an actual representative \(s_0\) from each nonempty
bracket, and apply \eqref{eq:gaussian-one-contaminant-bracket} with baseline
\(a_0=s_0^2p_0\) to each outer component at the same width. If the inner
bracket is \([L,U]\), replace \(L\) by \(L_+=\max\{L,0\}\), put
\(b=U-L_+\), and widen each component bracket \([\ell_j,u_j]\) to
\([(\ell_j-b)_+,u_j+b]\). Lemma~\ref{lem:contaminant-calculus} then gives
genuine composed endpoints through \(\Phi\), of total width at most \(\zeta\).
The logarithm of the shell bracket count is therefore at most
\begin{equation}
\label{eq:gaussian-shell-entropy}
 CD(d+1)\{\log M+\log D+\log_+(r/\eta)\}.
\end{equation}

The shell construction leading to
\eqref{eq:gaussian-shell-entropy} covers every positive Hellinger radius. To
cover the innermost Hellinger ball by a single bracket, apply the coefficient
bound and the fixed Taylor envelope just described. These two bounds give a fixed
\(H_0\in L^2(p_0)\) such that the inner ratio \(A=a/p_0\) satisfies
\(|A-1|\le C\rho_GH_0\). The outer part \(B=(p_G-a)/p_0\) satisfies
\[
 B/\rho_G\le C\sup_{|\theta|\le M}q_\theta/p_0.
\]
Completing the square shows that the supremum on the right has
\(L^2(p_0)\)-norm at most \(e^{C(1+M^2)}\). Since
\(|\sqrt{A+B}-1|\le|A+B-1|\), there is a common envelope \(F_M\) such that
\[
 |\sqrt{p_G/p_0}-1|\le\rho_GF_M,
 \qquad \|F_M\|_2\le e^{C(1+M^2)}.
\]
Thus the ball \(\rho_G\le\eta/(2\|F_M\|_2)\) lies in one bracket of width
\(\eta\). There are only
\(O\{1+M^2+\log_+(\delta/\eta)\}\) dyadic shells above it. Taking the union
of their bracket collections adds at most
\[
 C\{\log M+\log_+(\delta/\eta)\}
\]
to \eqref{eq:gaussian-shell-entropy}. Because \(D\le M^{C_0}\), the term
\(D\log D\) is absorbed by \(CD\log M\), and the shell-union term is
absorbed by the asserted entropy bound.

For the remaining region \(\rho_G\ge\rho_0\), no localization is needed.
Pointwise derivative brackets for Gaussian kernels, combined with a simplex
grid for the weights, give
\[
 \log N_{[]}\bigl(\eta^2,
 \{p_G:|\supp(G)|\le D,\ \supp(G)\subset B(0,M)\},L^1\bigr)
 \le CD(d+1)\{\log M+\log(1/\eta)\}.
\]
Here is the construction. A location cell \(I\), centered at \(\theta_I\),
is bracketed by
\([(q_{\theta_I}-\operatorname{diam}(I)E_I)_+,
q_{\theta_I}+\operatorname{diam}(I)E_I]\), where
\(E_I=\sup_{\theta\in I}|\nabla q_\theta|\) has uniformly bounded
\(L^1\)-norm. Quantizing weights at mesh \(h=c\eta^2/D\), and writing
\(k_j=\lfloor w_j/h\rfloor\), gives mixture endpoints
\(\sum_jk_jh\ell_j\) and \(\sum_j(k_j+1)hu_j\). Their \(L^1\)-distance is
at most \(C\eta^2\), and the number of feasible integer weight vectors is
at most \((C/\eta^2)^D\).
Taking square roots turns these into Hellinger brackets of width \(\eta\).
On the present region, \(\delta\ge\rho_0\), so
\(\log(1/\eta)\le C+\log(\delta/\eta)\). Combining the local and global
regions proves the lemma.
\end{proof}

We now insert the entropy estimate into the localized likelihood inequality.
The resulting bound applies simultaneously to all sparse competitors and
will be compared with Proposition~\ref{prop:gaussian-finite-constructive}.

\begin{proposition}
\label{prop:gaussian-finite-upper}
Fix \(C_0<\infty\). Let \(D_n\) be positive integers satisfying
\(D_n\le(\log n)^{C_0}\) and
\(D_n(d+1)\log\log n=o(n)\). Then
\[
    \mathbb P\left\{
    \sup_{|\supp(G)|\le D_n}L_n(G;G_0)
    \le C_{G_0,d,C_0}D_n(d+1)\log\log n
    \right\}\to1.
\]
\end{proposition}

\begin{proof}
Since \(G_0\) is fixed and finite, \(\mathbb P\{\max_i|X_i|\le C_{G_0}\sqrt{\log n}\}\to1\). On this event, Lemma~\ref{lem:convex-support} restricts all \(D_n\)-atomic competitors to \(B(0,C_{G_0}\sqrt{\log n})\). Apply Lemma~\ref{lem:gaussian-finite-local-bracket} with \(M=C_{G_0}\sqrt{\log n}\) and growth exponent \(2C_0+1\). The bracketing integral is bounded by \(C\delta\sqrt{A_n}\) for \(A_n=C_{G_0,d,C_0}D_n(d+1)\log\log n\), because \(\int_0^1\sqrt{1+\log(1/t)}\,dt<\infty\). Lemma~\ref{lem:wong-shen-local} gives the desired upper bound.
\end{proof}

\begin{proof}[Proof of Theorem~\ref{thm:gaussian-finite-main}]
By Proposition~\ref{prop:gaussian-finite-constructive}, the unrestricted likelihood is at least \(cJ_{n,d}\) with probability tending to one, where \(J_{n,d}\asymp_{G_0,d}(\log n)^{d/2}\). Proposition~\ref{prop:gaussian-finite-upper} gives the upper bound \(C D_n(d+1)\log\log n\) for all \(D_n\)-atomic competitors. Taking
\[
    D_n=
    \left\lfloor
        c_{G_0,d}\frac{(\log n)^{d/2}}{\log\log n}
    \right\rfloor
\]
with \(c_{G_0,d}\) small enough, no \(D_n\)-atomic mixing distribution can
maximize the likelihood on the intersection of the two high-probability
events. This proves the theorem.
\end{proof}

\subsection{Small Variance Misspecification}
\label{sec:gaussian-small-variance-mismatch}

The one-dimensional phase-transition analysis in
\cite{varianceD1Companion}
shows that, for \(\eps_n=n^{-r}\), the one-atom probability is nontrivial
when \(0<r<1/4\) and vanishes when \(r\ge1/4\). The following proposition
quantifies the support on the smaller-mismatch side: if \(\eps_n\) is below
\(n^{-1/4}\) by a fixed polynomial factor, then the lower bound has the same
order as under correct specification.

\begin{proposition}
\label{prop:gaussian-small-variance-mismatch}
Fix \(\delta>0\). Let \(X_1,\dots,X_n\) be i.i.d.\
\(\cN(0,1-\eps_n)\), where
\(0\le\eps_n\le n^{-1/4-\delta}\), and fit unit-variance Gaussian
location mixtures. There is a constant \(c_\delta>0\) such that every
NPMLE \(\wh G_n\) satisfies
\[
    \mathbb P\left\{
        |\supp(\wh G_n)|
        \ge
        c_\delta\frac{\sqrt{\log n}}{\log\log n}
    \right\}\to1.
\]
\end{proposition}

\begin{proof}
The proof is the one-dimensional construction from
Sections~\ref{sec:gaussian-constructive-gain} and
\ref{sec:gaussian-localized-upper}, with two changes: the score array is
controlled by truncation rather than raw fourth moments, and the variance
mismatch introduces a small deterministic damping term. We give the estimates
that justify both changes.

If the proposition holds for one value of \(\delta\), it also holds for every
larger value. We may therefore assume \(0<\delta<1/4\). Write
\(X_i=\sigma_nZ_i\), where \(\sigma_n^2=1-\eps_n\) and
\(Z_1,\ldots,Z_n\) are i.i.d.\ standard normal. Put
\(\gamma=\delta/2\), choose
\[
    \theta_0=\frac12-\frac\gamma2,
    \qquad
    \theta_1=\frac12-\frac\gamma4,
\]
and let \(u_n\) satisfy \(n\mathbb P\{Z>u_n\}=1\) for
\(Z\sim\cN(0,1)\). A constant-spaced grid
\(\mathcal A_n=\{a_1,\dots,a_J\}\) in
\([\theta_0u_n,\theta_1u_n]\) has
\(J\asymp_\delta\sqrt{\log n}\). For
\[
    q_a(z)=e^{az-a^2/2},
    \qquad
    h_a(z)=\frac{q_a(z)-1}{(e^{a^2}-1)^{1/2}},
\]
direct Gaussian integration gives
\[
    \mathbb E h_a(Z)h_b(Z)
    =
    \frac{e^{ab}-1}{\sqrt{(e^{a^2}-1)(e^{b^2}-1)}}
    \le 2e^{-(a-b)^2/2}
    \qquad(a\ne b).
\]
Choosing the grid spacing large makes the off-diagonal Gram row sums smaller
than \(1/10\).

The earlier construction used a smaller score radius and controlled the
empirical Gram matrix by raw fourth moments. Here the grid approaches
\(a=u_n/2\), where those moments are too large, so we replace that step by
the following truncation estimates.
The strict inequality \(\theta_1<1/2\) gives \(c>0\) and
\(\rho_n=n^{-c}\) such that, uniformly over \(a,b\) in the grid,
\begin{align*}
    \mathbb E\!\left[h_a(Z)^2
        \mathbf 1_{\{|h_a(Z)|>\rho_n\sqrt n\}}\right]
    +n\mathbb P\{|h_a(Z)|>\rho_n\sqrt n\}
    &\le \rho_n,\\
    \mathbb E\!\left[|h_a(Z)h_b(Z)|
        \mathbf 1_{\{|h_a(Z)h_b(Z)|>n\rho_n\}}\right]
    &\le \rho_n,\\
    \mathbb E\!\left[(h_a(Z)h_b(Z))^2
        \mathbf 1_{\{|h_a(Z)h_b(Z)|\le n\rho_n\}}\right]
    &\le n\rho_n.
\end{align*}
To verify these estimates, note that
\((e^{a^2}-1)^{1/2}\asymp e^{a^2/2}\). If \(a=\theta u_n\), the first
tail event forces
\[
    Z\ge
    \left(\theta+\frac1{4\theta}-\frac{c}{2\theta}+o(1)\right)u_n.
\]
At \(c=0\), the coefficient is uniformly larger than both \(1\) and
\(2\theta\) on \([\theta_0,\theta_1]\). Mills' bound controls the ordinary
tail, while
\[
    h_a(z)^2\phi_1(z)
    \le C\phi_1(z-2a)+Ce^{-a^2}\phi_1(z)
\]
controls its second moment under the tilted normal law. The product estimates
then follow from Cauchy--Schwarz, after reducing \(c\). More explicitly,
repeating this tilted-tail calculation at the larger threshold
\(\sqrt{n\rho_n}\) first gives, for
some \(c_0>0\),
\[
 \sup_a\mathbb E\!\left[h_a(Z)^2
   \mathbf 1_{\{|h_a(Z)|>\sqrt{n\rho_n}\}}\right]
 \le n^{-c_0}.
\]
Choose \(c<c_0/2\). Since
\(\{|h_ah_b|>n\rho_n\}\) is contained in the union of the two events
\(\{|h_a|>\sqrt{n\rho_n}\}\) and
\(\{|h_b|>\sqrt{n\rho_n}\}\), Cauchy--Schwarz gives the product-tail
bound above. The final truncated second-moment bound follows from
\((h_ah_b)^2\mathbf 1_{\{|h_ah_b|\le n\rho_n\}}
\le n\rho_n|h_ah_b|\).

Now set
\[
    d_j=\frac1{\sqrt n}\sum_{i=1}^nh_{a_j}(Z_i),
    \qquad
    \wh H_{jk}=\frac1n\sum_{i=1}^nh_{a_j}(Z_i)h_{a_k}(Z_i),
    \qquad
    S=\sum_{j=1}^J(d_j^+)^2.
\]
Truncating the products at \(n\rho_n\), applying Chebyshev's inequality to
each entry, and taking a union bound gives
\(\|\wh H-H\|_{\op}=o_{\mathbb P}(1)\). Truncating the individual scores at
\(\rho_n\sqrt n\), and repeating the bounded positive-score argument from
Proposition~\ref{prop:gaussian-finite-constructive}, gives
\[
    S=O_{\mathbb P}(J),
    \qquad
    \mathbb P\{S\ge cJ\}\to1.
\]
Indeed, \(J\) is polylogarithmic whereas \(\rho_n\) is a negative power of
\(n\), so the entrywise union bounds, the operator-norm conversion, and the
one- and two-coordinate Berry--Esseen errors all vanish. More explicitly,
with \(T_n=\rho_n\sqrt n\),
\[
 \left|\mathbb E h_a(Z)\mathbf 1_{\{|h_a(Z)|>T_n\}}\right|
 \le
 \sqrt{\mathbb P\{|h_a(Z)|>T_n\}
       \mathbb E[h_a(Z)^2\mathbf 1_{\{|h_a(Z)|>T_n\}}]}
 \le \frac{\rho_n}{\sqrt n}.
\]
Thus recentering moves a normalized score by at most \(\rho_n\), and
variance normalization changes it by \(O(\rho_n)\). The probability that
any of the \(nJ\) summands is truncated is at most \(J\rho_n=o(1)\), so the
Gaussian approximation transfers back to the original scores \(d_j\).
Finally, choose \(\zeta>0\) so that
\(2\theta_1(1+\zeta-\theta_1)<1/2\). With probability tending to one,
\(\max_i|Z_i|\le(1+\zeta)u_n\), and the explicit score formula then gives,
for some \(c_1>0\),
\[
    \max_{i,j}|h_{a_j}(Z_i)|\le n^{1/2-c_1}
\]
with probability tending to one. Thus both ingredients used in the correctly
specified construction remain valid: the empirical score vector has the same
Gaussian approximation, and every observed score is uniformly small enough
for the likelihood Taylor expansion.

We now transfer the correctly specified construction from the standardized
sample \(Z_1,\ldots,Z_n\) to the observed sample. Put
\(\kappa_n=\eps_n/\{2(1-\eps_n)\}\). A fitted atom
\(b_j=a_j/\sigma_n\) has density ratio
\[
    \frac{\phi_1(X_i-b_j)}{\phi_1(X_i)}
    =e^{-\kappa_na_j^2}q_{a_j}(Z_i).
\]
Thus \(e^{-\kappa_na_j^2}=1+O(\eps_n\log n)=1+o(1)\) uniformly on the
grid. As in Proposition~\ref{prop:gaussian-finite-constructive}, fix
\(\eta=1/5\), set
\[
    r_j=\frac{d_j^+}{1+\eta},
    \qquad
    t_j=\frac{r_j}{\sqrt n\,(e^{a_j^2}-1)^{1/2}},
\]
and move mass \(t_j\) from zero to \(b_j\). Since
\(\sum_jr_j=O_{\mathbb P}(J)\), the total moved mass is
\(o_{\mathbb P}(1)\).

Relative to the correctly specified calculation, the only additional
constant in each fitted density ratio is
\[
    c_n=\frac1{\sqrt n}\sum_{j=1}^J
        \frac{r_j\{e^{-\kappa_na_j^2}-1\}}
             {(e^{a_j^2}-1)^{1/2}}.
\]
Its contribution to the linear likelihood term satisfies
\[
    n|c_n|
    \le O_{\mathbb P}(J)(\log n)
        n^{1/2-(1/4+\delta)-\theta_0^2+o(1)}
    =o_{\mathbb P}(J),
\]
because \(\theta_0^2>1/4-\delta\). The damping factors also change the
linear and quadratic score terms by only \(o_{\mathbb P}(J)\). The constant
\(c_n\) creates the additional terms
\(2c_n\sum_i v_i\) and \(nc_n^2\) in the quadratic expansion, where
\[
 v_i=\frac1{\sqrt n}\sum_j
 r_je^{-\kappa_na_j^2}h_{a_j}(Z_i),
 \qquad
 \sum_i v_i=O_{\mathbb P}(J).
\]
They are \(o_{\mathbb P}(J)\) because
\(n|c_n|=o_{\mathbb P}(J)\) and \(J=o(n)\). The bound
\(\max_{i,j}|h_{a_j}(Z_i)|\le n^{1/2-c_1}\) makes the Taylor remainder
\(o_{\mathbb P}(J)\), exactly
as in Proposition~\ref{prop:gaussian-finite-constructive}. The quadratic
likelihood expansion from that proposition now applies with the damping,
centering, and Taylor errors just bounded. It gives a feasible mixing
distribution with
likelihood gain at least \(cJ\asymp_\delta\sqrt{\log n}\) with probability
tending to one.

The sparse upper bound transfers without any new entropy calculation. For any
fitted atom \(b\), put \(t=\sigma_nb\). Then
\[
    \frac{\phi_1(X_i-b)}{\phi_1(X_i)}
    =e^{-\kappa_nt^2}q_t(Z_i)\le q_t(Z_i).
\]
After mixing, every \(D\)-atomic likelihood for \(X_1,\dots,X_n\) is therefore
bounded by a \(D\)-atomic correctly specified likelihood for
\(Z_1,\dots,Z_n\). Proposition~\ref{prop:gaussian-finite-upper} bounds the
correctly specified likelihood by \(CD\log\log n\) with probability tending
to one. Comparing the two
likelihood scales and taking
\(D=c_\delta\sqrt{\log n}/\log\log n\), with \(c_\delta\) small enough,
proves the proposition.
\end{proof}

The one-dimensional restriction reflects the mismatch scale on which the
bound is informative, not a limitation of the score construction. The same
construction gives the following higher-dimensional statement.

\begin{remark}
For every fixed \(d\ge2\), replacing the one-dimensional grid
\(\mathcal A_n\) by a
constant-separated grid in the annulus
\[
    \{a\in\R^d:\theta_0u_n\le |a|\le\theta_1u_n\}
\]
gives \(J\asymp_{\delta,d}(\log n)^{d/2}\) score locations. For the
multivariate Gaussian scores, the exact covariance formula implies
\[
 |\mathbb E h_a(Z)h_b(Z)|
 \le C\left\{e^{-|a-b|^2/2}
       +e^{-(|a|^2+|b|^2)/2}\right\}.
\]
Choosing the separation constant large makes the off-diagonal row sums of
the first term arbitrarily small; the second has row sum at most
\(J e^{-\theta_0^2u_n^2}=o(1)\). Thus the truncation, damping, and
sparse-likelihood arguments are otherwise unchanged.
Repeating the score and likelihood calculations on this annular grid gives
\[
    |\supp(\wh G_n)|
    \ge
    c_{\delta,d}\frac{(\log n)^{d/2}}{\log\log n}
\]
with probability tending to one whenever
\(\eps_n\le n^{-1/4-\delta}\).

We state Proposition~\ref{prop:gaussian-small-variance-mismatch} only in one
dimension because this polynomial condition is tied to the correct
one-dimensional boundary. In higher dimensions it is far below the actual
transition scale: the companion paper finds
\(\eps_n\asymp1/\log n\) for \(d=2\) and
\(\eps_n\asymp(\log\log n)/\log n\) for \(d\ge3\)
\cite{varianceMultivariateCompanion}. Proving quantitative support growth
near those scales would require uniform control over the continuum of angular
directions near the radial extremes, not merely the sub-edge Lindeberg
argument used here.
\end{remark}


\section{Poisson Mixture Analogue}
\label{sec:poisson-analogue}

This section proves Theorem~\ref{thm:poisson-finite-main}. As in the Gaussian
case, Proposition~\ref{prop:poisson-finite-constructive} constructs a large
unrestricted likelihood gain and Proposition~\ref{prop:poisson-finite-upper}
bounds every \(D\)-atomic competitor. The Poisson proof differs in its score
analysis: truncation at the observed sample scale replaces the Gaussian
moment estimates and permits perturbations across the full exposed tail
window.

For the Poisson theorem, write
\[
    G_0=\sum_{m=1}^K\pi_m\delta_{\lambda_m},
    \qquad
    0\le \lambda_1<\cdots<\lambda_K=\lambda_*,
    \qquad \lambda_*>0,
\]
and write \(p_0=p_{G_0}\). The relative log-likelihood is
\[
    L_n^{\mathrm{Pois}}(G;G_0)
    =
    \sum_{i=1}^n\log\frac{p_G(X_i)}{p_0(X_i)}.
\]
Write \(\pi_*=\pi_K\). The construction perturbs the component at
\(\lambda_*\). This component is identifiable from the upper tail because,
among the finitely many true Poisson laws, it dominates the probability of
large counts.

As in the Gaussian case, a sample-dependent compactness reduction controls sparse competitors. For Poisson mixtures, the compact set is the interval up to the sample maximum.

\begin{lemma}
\label{lem:poisson-support}
Let \(M_n=\max_{1\le i\le n}X_i\). For every probability measure \(G\) on \([0,\infty)\), there exists a probability measure \(\widetilde G\) supported on \([0,M_n]\) such that \(p_{\widetilde G}(X_i)\ge p_G(X_i)\) for every \(1\le i\le n\). Moreover, if \(G\) has at most \(D\) atoms, then so does \(\widetilde G\).
\end{lemma}

\begin{proof}
Let \(T(\theta)=\min(\theta,M_n)\) and set \(\widetilde G=T_\#G\). Fix an observed value \(k\le M_n\). If \(M_n=0\), then \(k=0\) and \(p_\theta(0)=e^{-\theta}\le1=p_{\theta=0}(0)\). Otherwise, \(d\log p_\theta(k)/d\theta=k/\theta-1\), so \(\theta\mapsto p_\theta(k)\) is decreasing on \([M_n,\infty)\). Hence \(p_{T(\theta)}(k)\ge p_\theta(k)\) for every \(\theta\ge0\), and integration over \(G\) proves the lemma.
\end{proof}

\subsection{Scores Near the Largest Mean}
\label{sec:poisson-largest-mean-scores}

The upper tail exposes the largest true mean \(\lambda_*\). The score estimates
below show that separated outward perturbations of this component are nearly
orthogonal after truncation; these are the inputs to
Proposition~\ref{prop:poisson-finite-constructive}.

For \(a\ge0\), set \(\theta_a=\lambda_*+a\sqrt{\lambda_*}\) and
\[
    \Delta_a(k)=p_{\theta_a}(k)-p_{\lambda_*}(k),
    \qquad
    g_a(k)=\frac{\Delta_a(k)}{p_0(k)},
\]
\[
    \sigma_a^2=\sum_{k=0}^\infty\frac{\Delta_a(k)^2}{p_0(k)},
    \qquad
    h_a(k)=\frac{g_a(k)}{\sigma_a}.
\]
Then \(\E_{p_0}h_a(X)=0\) and \(\E_{p_0}h_a(X)^2=1\). Also
\begin{equation}
\label{eq:poisson-ratio-comparison}
    \pi_*
    \le
    \frac{p_0(k)}{p_{\lambda_*}(k)}
    =
    \pi_*+
    \sum_{m<K}\pi_m e^{\lambda_* -\lambda_m}
    \left(\frac{\lambda_m}{\lambda_*}\right)^k
    \le C_{G_0},
    \qquad k\ge0.
\end{equation}
In \eqref{eq:poisson-ratio-comparison}, the term corresponding to
\(\lambda_m=0\) at \(k=0\) is interpreted using \(0^0=1\).
Thus replacing the full denominator \(p_0\) by the single-component
denominator \(p_{\lambda_*}\) changes a score by at most constant factors.

The next lemma gives the covariance and concentration estimates needed to
combine many such scores. A direct moment bound would be governed by counts
far beyond the observed sample range. The truncated estimates discard that
irrelevant part of the likelihood ratio and remain uniform up to
\(a^2\asymp\log n/\log\log n\).

\begin{lemma}
\label{lem:poisson-finite-score-truncation}
Fix \(\gamma\in(0,1)\), and put \(R_n=\sqrt{\gamma\log n/\log\log n}\). There are constants \(A_0=A_0(G_0)\), \(c=c(G_0,\gamma)>0\), and \(C=C(G_0)<\infty\) such that, uniformly for \(A_0\le a,b\le R_n\),
\[
    C^{-1}e^{a^2}\le \sigma_a^2\le C e^{a^2},
    \qquad
    |\E_{p_0}h_a(X)h_b(X)|\le C e^{-(a-b)^2/2}.
\]
Moreover, with \(\rho_n=e^{-c\log n/\log\log n}\), uniformly for \(A_0\le a,b\le R_n\),
\[
    \E\bigl[h_a(X)^2\mathbf 1_{\{|h_a(X)|>\rho_n\sqrt n\}}\bigr]
    \le \rho_n,
    \qquad
    n\mathbb P\{|h_a(X)|>\rho_n\sqrt n\}
    \le \rho_n,
\]
\[
    \E\bigl[
        |h_a(X)h_b(X)|
        \mathbf 1_{\{|h_a(X)h_b(X)|>n\rho_n\}}
    \bigr]
    \le \rho_n,
\]
and
\[
    \E\bigl[
        (h_a(X)h_b(X))^2
        \mathbf 1_{\{|h_a(X)h_b(X)|\le n\rho_n\}}
    \bigr]
    \le n\rho_n.
\]
\end{lemma}

\begin{proof}
Let \(\tau_a=(e^{a^2}-1)^{1/2}\) and let
\[
    h_{a,*}(k)=
    \frac{p_{\theta_a}(k)/p_{\lambda_*}(k)-1}{\tau_a}
\]
be the one-component score at \(\lambda_*\). By
\eqref{eq:poisson-ratio-comparison}, \(\sigma_a\asymp\tau_a\) and
\(|h_a(k)|\asymp |h_{a,*}(k)|\), uniformly for \(a\ge A_0\) and
\(k\ge0\), after taking \(A_0\) fixed and large enough. Under the pure
\(p_{\lambda_*}\) law, the unnormalized score covariance is exactly
\(e^{ab}-1\). To compare this formula with the covariance under the full
mixture, write
\[
    W(k)=\frac{p_0(k)}{p_{\lambda_*}(k)}
    =\pi_*+\sum_{m<K}c_m r_m^k,
    \qquad 0\le r_m<1.
\]
Thus \(W(k)^{-1}=\pi_*^{-1}+O(r^k)\) for some fixed \(r<1\). If
\(q_a=p_{\theta_a}/p_{\lambda_*}\), then
\[
    \E_{\lambda_*}[q_aq_b]=e^{ab},
\]
whereas
\[
    \E_{\lambda_*}[r^Xq_aq_b]
    =
    \exp\!\left\{
        rab-(1-r)\bigl((a+b)\sqrt{\lambda_*}+\lambda_*\bigr)
    \right\}.
\]
The analogous formulas with \(q_aq_b\) replaced by \(q_a\), \(q_b\), or
\(1\) show that
\[
    \sum_k\frac{\Delta_a(k)\Delta_b(k)}{p_0(k)}
    =
    \pi_*^{-1}(e^{ab}-1)
    +O\!\left(
        e^{rab-c(a+b)}+e^{-ca}+e^{-cb}+1
    \right).
\]
For \(a=b\), this proves \(\sigma_a^2\asymp e^{a^2}\). For arbitrary
\(a,b\ge A_0\), division by
\(\sigma_a\sigma_b\asymp e^{(a^2+b^2)/2}\) gives
\[
    |\E h_a(X)h_b(X)|
    \le C e^{-(a-b)^2/2};
\]
the error term is smaller by the additional factor
\(e^{-(1-r)ab-c(a+b)}\).

We next locate the counts at which the score exceeds the truncation level.
Fix \(C_0\) so that \(|h_a(k)|\le C_0|h_{a,*}(k)|\) throughout the stated
range, and put \(T_n=C_0^{-1}\rho_n\sqrt n\). Since
\[
 \frac{p_{\theta_a}(k)}{p_{\lambda_*}(k)}
 =e^{-a\sqrt{\lambda_*}}
  \left(1+\frac a{\sqrt{\lambda_*}}\right)^k,
\]
the event \(|h_a(X)|>\rho_n\sqrt n\) forces \(X\ge k_a\), where
\[
    k_a
    =
    \left\lceil
    \frac{
        a\sqrt{\lambda_*}+\log(1+T_n\tau_a)
    }
    {\log(1+a/\sqrt{\lambda_*})}
    \right\rceil .
\]
Equivalently,
\[
 k_a=
 \frac{
        \frac12\log n+\frac12a^2-c\log n/\log\log n+O_{G_0}(a+1)
      }
      {\log(1+a/\sqrt{\lambda_*})}
 +O(1).
\]
The final \(O(1)\) is the rounding error. After multiplication by
\(\log(1+a/\sqrt{\lambda_*})=O(\log\log n)\), it is negligible compared
with the margins of order \(\log n/\log\log n\) below.
Uniformly for \(A_0\le a\le R_n\), and with \(c>0\) small enough, there are constants \(\delta_0,c_1,c_2>0\), depending only on \(G_0\) and \(\gamma\), such that this threshold satisfies
\[
    k_a\ge(1+\delta_0)(\sqrt{\lambda_*}+a)^2,
    \qquad
    k_a\log\frac{k_a}{e\lambda_*}
    \ge
    \log n+c_1\frac{\log n}{\log\log n},
\]
and
\[
    k_a\log\frac{k_a}{e(\sqrt{\lambda_*}+a)^2}
    +(\sqrt{\lambda_*}+a)^2
    \ge
    c_2\frac{\log n}{\log\log n}.
\]
To verify these inequalities uniformly in \(a\), write
\(L=\log n\), \(\ell=\log L\), and
\(y=a^2\ell/L\). On every fixed interval \(y\in[\delta,\gamma]\),
\[
    k_a
    =
    \frac{L}{\ell}
    \left[
        1+
        \frac{\log\ell-\log y+\log\lambda_*+y-2c+o(1)}{\ell}
    \right].
\]
Consequently, if \(I_\mu(k)=k\log(k/\mu)-k+\mu\) denotes the Poisson
Chernoff exponent, then, uniformly on this interval,
\[
    I_{\lambda_*}(k_a)
    =
    L+
    \{y-\log y-1-2c+o(1)\}\frac{L}{\ell},
\]
and, for \(\mu_a=(\sqrt{\lambda_*}+a)^2\),
\[
    I_{\mu_a}(k_a)
    =
    \{y-\log y-1+o(1)\}\frac{L}{\ell}.
\]
The function \(y-\log y-1\) is decreasing on \((0,1)\), so the two rate
functions have a fixed positive margin on \([\delta,\gamma]\). Also
\(k_a/\mu_a=1/y+o(1)\) there, uniformly, and is therefore bounded away
from one.

The expansion is not uniform when \(y\) approaches zero. For that range, fix
\(\delta\in(0,\min\{\gamma,e^{-3}\})\), and write
\[
 w_a=2\log(1+a/\sqrt{\lambda_*})
     =\log(\mu_a/\lambda_*).
\]
Since \(a\le R_n\), the \(O(a+1)\) term in the definition of \(k_a\) is
\(o(L/\ell)\), uniformly. Hence
\begin{equation}
\label{eq:poisson-threshold-product}
 k_aw_a\ge L+\mu_a-3cL/\ell
\end{equation}
for all sufficiently large \(n\). If \(w_a\le\ell/2\), then
\(\mu_a=\lambda_*e^{w_a}\le\lambda_*\sqrt L\), while
\(k_a\ge(2-o(1))L/\ell\). Thus \(k_a/\mu_a\to\infty\) uniformly in this
range and, eventually,
\[
 I_{\mu_a}(k_a)\ge k_a\ge L/\ell.
\]

Suppose instead that \(w_a>\ell/2\), and put
\(z=\mu_a\ell/L\). The assumption \(y\le\delta\) gives
\(z\le2\delta\) for large \(n\), while
\(w_a=\ell-\log\ell+\log(z/\lambda_*)\le\ell\). If
\(r_a=k_a/\mu_a\), then \eqref{eq:poisson-threshold-product} yields,
uniformly in this range,
\[
 r_a\ge\frac{\ell+z-3c}{zw_a}
     \ge\frac{1-o(1)}{z}.
\]
Choose a fixed \(\varepsilon\in(0,1/4)\). For large \(n\), the right side
is at least \((1-\varepsilon)/z\). Since
\(f(r)=r\log r-r+1\) is increasing for \(r>1\),
\[
 \frac{\ell}{L}I_{\mu_a}(k_a)
 =z f(r_a)
 \ge (1-\varepsilon)\log\frac{1-\varepsilon}{z}
      -(1-\varepsilon)+z.
\]
The expression on the right is decreasing for
\(0<z<1-\varepsilon\); at \(z=2\delta\) it is a positive constant
\(m_\delta\). Consequently
\(I_{\mu_a}(k_a)\ge m_\delta L/\ell\), and
\(k_a/\mu_a\) is uniformly bounded away from one, throughout
this second range. Put \(\bar m_\delta=\min\{1,m_\delta\}\). The two
\(w_a\)-ranges therefore give
\(I_{\mu_a}(k_a)\ge\bar m_\delta L/\ell\) throughout
\(y\le\delta\). Finally, the exact identity
\[
 I_{\lambda_*}(k_a)-I_{\mu_a}(k_a)
 =k_aw_a+\lambda_*-\mu_a
\]
and \eqref{eq:poisson-threshold-product} give
\[
 k_a\log\frac{k_a}{e\lambda_*}
 =I_{\lambda_*}(k_a)-\lambda_*
 \ge L+(\bar m_\delta-3c)L/\ell.
\]
Choose \(c>0\) smaller than one quarter of
\(\gamma-\log\gamma-1\) and smaller than \(\bar m_\delta/6\).
The compact and small-\(y\) ranges together prove all three threshold
inequalities, after decreasing \(\delta_0,c_1,c_2\).

The first rate function controls the threshold probability under
\(p_{\lambda_*}\), and the second controls the truncated second moment after
exponential tilting. The identity
\[
    \left(\frac{p_{\theta_a}(k)}{p_{\lambda_*}(k)}\right)^2p_{\lambda_*}(k)
    =
    e^{a^2}p_{(\sqrt{\lambda_*}+a)^2}(k)
\]
therefore gives, by Chernoff's bound,
\[
    \E_{p_{\lambda_*}}
    \bigl[h_{a,*}(X)^2\mathbf 1_{\{|h_{a,*}(X)|>\rho_n\sqrt n\}}\bigr]
    \le e^{-c\log n/\log\log n}
\]
and
\[
    n\mathbb P_{p_{\lambda_*}}
    \{|h_{a,*}(X)|>\rho_n\sqrt n\}
    \le e^{-c\log n/\log\log n}.
\]
For the product tail, use
\(
\{|h_a h_b|>n\rho_n\}
\subset
\{|h_a|>\sqrt{n\rho_n}\}\cup\{|h_b|>\sqrt{n\rho_n}\}
\).
Repeating the one-coordinate calculation at the larger threshold
\(\sqrt{n\rho_n}\) gives, for some \(c_3>0\), a uniform second-moment
tail bound \(\exp\{-c_3\log n/\log\log n\}\). Cauchy--Schwarz on the two
events therefore bounds the product tail by
\(2\exp\{-c_3\log n/(2\log\log n)\}\). Decreasing the constant \(c\) in
the definition of \(\rho_n\) makes this at most \(\rho_n\). Finally,
\[
    (h_a h_b)^2\mathbf 1_{\{|h_a h_b|\le n\rho_n\}}
    \le n\rho_n |h_a h_b|,
\]
and \(\E|h_a h_b|\le1\).
\end{proof}

\subsection{Constructive Likelihood Gain}
\label{sec:poisson-constructive-gain}

The following proposition constructs the unrestricted likelihood gain. It
moves a small amount of mass from \(\lambda_*\) to a separated grid of larger
means \(\lambda_*+a_j\sqrt{\lambda_*}\).

\begin{proposition}
\label{prop:poisson-finite-constructive}
Fix \(\gamma\in(0,1)\), and put \(R_n=\sqrt{\gamma\log n/\log\log n}\). Choose \(\Delta\) large enough, let \(a_j=A_0+\Delta(j-1)\le R_n\), and let \(J\) be the number of such points. Then \(J\asymp_{G_0,\gamma,\Delta}\sqrt{\log n/\log\log n}\), and there is a random mixing distribution \(G_n\) such that
\[
    \mathbb P\left\{
        L_n^{\mathrm{Pois}}(G_n;G_0)
        \ge c_{G_0,\gamma,\Delta}J
    \right\}\to1.
\]
\end{proposition}

\begin{proof}
For \(1\le j\le J\), write \(h_j=h_{a_j}\),
\[
    d_j=\frac1{\sqrt n}\sum_{i=1}^n h_j(X_i),
    \qquad
    H_{jl}=\E h_j(X)h_l(X),
    \qquad
    \wh H_{jl}=\frac1n\sum_{i=1}^n h_j(X_i)h_l(X_i).
\]
Thus \(d_j\) is the empirical score in direction \(a_j\), while \(H\) and
\(\wh H\) are the population and empirical Gram matrices of these scores.
Choose \(\Delta\) so large that the off-diagonal row sums of \(H\) are at most \(1/10\); this is possible by Lemma~\ref{lem:poisson-finite-score-truncation}. The truncation bounds in Lemma~\ref{lem:poisson-finite-score-truncation}, followed by Chebyshev's inequality and a union bound over \(j,l\), give
\[
    \|\wh H-H\|_{\op}=o_{\mathbb P}(1).
\]
Here are the details. Put
\[
    Y_i^{jl}
    =
    h_j(X_i)h_l(X_i)
    \mathbf 1_{\{|h_j(X_i)h_l(X_i)|\le n\rho_n\}}.
\]
The truncation lemma gives
\[
    |\E Y_i^{jl}-H_{jl}|\le\rho_n,
    \qquad
    \Var\left(\frac1n\sum_{i=1}^nY_i^{jl}\right)\le\rho_n.
\]
Moreover, with probability tending to one,
\(\max_{i,j}|h_j(X_i)|\le\rho_n\sqrt n\), because
\(nJ\max_j\mathbb P\{|h_j(X)|>\rho_n\sqrt n\}\le J\rho_n\to0\).
On this event no product is truncated. Chebyshev's inequality, followed by a
union bound over the \(J^2\) pairs, therefore yields
\[
    \max_{j,l}|\wh H_{jl}-H_{jl}|
    =O_{\mathbb P}(\rho_n^{1/4}).
\]
Since \(J\rho_n^{1/4}\to0\), the operator-norm claim follows.

Let \(S=\sum_{j=1}^J(d_j^+)^2\). We claim that a positive fraction of
the score coordinates contribute to \(S\). Fix a large constant \(B\) and
let \(\psi_B(x)=(x^+)^2\wedge B\). To obtain a uniform Gaussian
approximation, truncate each \(h_j\) at \(\rho_n\sqrt n\), then recenter and
renormalize it. If
\[
 A_{jl}=\{|h_j|>\rho_n\sqrt n\}\cup
        \{|h_l|>\rho_n\sqrt n\},
\]
then the product-tail estimate gives
\[
 \mathbb E|h_jh_l|\mathbf 1_{A_{jl}}
 \le \rho_n+n\rho_n\,\mathbb P(A_{jl})
 \le \rho_n+2\rho_n^2.
\]
Also \(n\mathbb P(A_{jl})\le2\rho_n\), so the original pair of normalized
sums agrees with its truncated version with probability at least
\(1-2\rho_n\).
Thus truncation changes each one- or two-coordinate covariance matrix by
\(O(\rho_n)\), and the matrices remain uniformly nonsingular. The truncated
third absolute moment is at most
\(\rho_n\sqrt n\,\mathbb E h_j(X)^2\). The fixed-dimensional
Berry--Esseen bound in Wasserstein distance therefore gives error
\(O(\rho_n)\) for every bounded Lipschitz function of one or two normalized
coordinates. This applies to \(\psi_B\) and to the product test function
\[
 (x,y)\longmapsto\psi_B(x)\psi_B(y).
\]
Because \(\rho_n\) decays faster than any power of \(J\), the total
approximation error over all \(J^2\) pairs is \(o(1)\). Cauchy--Schwarz also
gives
\[
 \left|\mathbb E h_j(X)\mathbf 1_{\{|h_j(X)|>\rho_n\sqrt n\}}\right|
 \le \frac{\rho_n}{\sqrt n}.
\]
Thus recentering shifts a normalized sum by at most \(\rho_n\), and variance
normalization changes it by \(O(\rho_n)\). This completes the transfer from
the truncated array back to the original coordinates \(d_j\).

For a centered Gaussian pair
\((Z_j,Z_l)\) with correlation \(H_{jl}\), the Gaussian maximal
correlation inequality gives
\[
    |\Cov(\psi_B(Z_j),\psi_B(Z_l))|
    \le C_B|H_{jl}|.
\]
The uniform Gaussian approximation and the row-sum bound for \(H\) imply
that \(\sum_j\psi_B(d_j)\) has variance
\(O_B(J)+o(J^2)\). Choose \(B\) so that
\(\E\psi_B(Z)>1/3\) for \(Z\sim\mathcal N(0,1)\). Its mean is then at
least \(J/3+o(J)\). Since \(\sum_j\psi_B(d_j)\le S\), Chebyshev's
inequality shows that \(S\ge cJ\) with high probability. The upper bound
follows from \(\E\sum_jd_j^2=J\). Thus
\[
    \mathbb P\{S\ge cJ\}\to1,
    \qquad
    S=O_{\mathbb P}(J).
\]

Fix \(\eta=1/5\), let \(r_j=d_j^+/(1+\eta)\), and set
\(t_j=r_j/(\sqrt n\,\sigma_{a_j})\). Since
\(\inf_j\sigma_{a_j}>0\) and
\(\sum_jr_j\le\sqrt{JS}=O_{\mathbb P}(J)=o_{\mathbb P}(\sqrt n)\),
we have \(\sum_jt_j=o_{\mathbb P}(1)\). Hence
\(\sum_jt_j\le\pi_*\) with probability tending to one. On this event define
\[
    G_n=
    G_0-
    \Bigl(\sum_{j=1}^Jt_j\Bigr)\delta_{\lambda_*}
    +
    \sum_{j=1}^Jt_j\delta_{\lambda_*+a_j\sqrt{\lambda_*}},
\]
and set \(G_n=G_0\) otherwise.

On the event \(\sum_jt_j\le\pi_*\), for each observation,
\[
    \frac{p_{G_n}(X_i)}{p_0(X_i)}
    =
    1+u_i,
    \qquad
    u_i=\frac1{\sqrt n}\sum_{j=1}^Jr_jh_j(X_i).
\]
This feasibility event has probability tending to one. Intersect it with the
event, also of probability tending to one, on which
\(\max_{i,j}|h_j(X_i)|\le\rho_n\sqrt n\), supplied by
Lemma~\ref{lem:poisson-finite-score-truncation} and a union bound. On their
intersection, \(\max_i|u_i|\le\rho_n\sum_jr_j=o_{\mathbb P}(1)\), and Taylor
expansion gives
\[
    L_n^{\mathrm{Pois}}(G_n;G_0)
    =
    r^\top d-\frac12 r^\top\wh H r+o_{\mathbb P}(J).
\]
Indeed, \(\sum_i u_i^2=r^\top\wh H r=O_{\mathbb P}(J)\), so the sum of
the cubic remainders is bounded by
\(O(\max_i|u_i|\sum_i u_i^2)=o_{\mathbb P}(J)\).
Since \(r^\top d=S/(1+\eta)\) and
\(r^\top\wh H r\le(1.2+o_{\mathbb P}(1))S/(1+\eta)^2\), the coefficient
of \(S\) is positive. The lower bound on \(S\) completes the proof.
\end{proof}

\subsection{Localized Upper Bound}
\label{sec:poisson-localized-upper}

We now bound every \(D\)-atomic Poisson competitor. Lemma
~\ref{lem:poisson-support} restricts its means to the interval ending at the
sample maximum, which is of order \(\log n/\log\log n\). As in the Gaussian
case, the required entropy estimate must depend on this parameter radius only
through its logarithm.

For \(D\ge1\), \(M\ge2\), and \(0<\delta\le\sqrt2\), define
\[
    \mathcal F_D^{\mathrm{Pois},G_0}(M,\delta)
    =
    \left\{
        \sqrt{p_G/p_0}-1:
        |\supp(G)|\le D,
        \ \supp(G)\subset[0,M],
        \ \|\sqrt{p_G/p_0}-1\|_{L^2(p_0)}\le\delta
    \right\}.
\]

The coordinate \(u=\sqrt\theta\) makes this logarithmic dependence
available: square-root Poisson kernels have Gaussian affinity
\(\sum_k\sqrt{p_{u^2}(k)p_{v^2}(k)}=e^{-(u-v)^2/2}\). The next lemma
combines this geometry with local identifiability around the true finite
mixture.

\begin{lemma}
\label{lem:poisson-finite-local}
For fixed \(G_0\) and fixed \(C_0<\infty\), there exists
\(C_{G_0,C_0}<\infty\) such that, for every \(M\ge2\),
\(1\le D\le M^{C_0}\), and \(0<\eta\le\delta\le\sqrt2\),
\[
    \log N_{[]}
    \bigl(
        \eta,
        \mathcal F_D^{\mathrm{Pois},G_0}(M,\delta),
        L^2(p_0)
    \bigr)
    \le
        C_{G_0,C_0}D\bigl(\log M+\log(\delta/\eta)\bigr).
\]
\end{lemma}

\begin{proof}
Put \(\rho_G=\|\sqrt{p_G/p_0}-1\|_2\). We first record the local geometry
needed when \(\rho_G\le\rho_0\). Choose relatively open, bounded intervals
\(A_i\) around the true means \(\lambda_i\), with disjoint closures, and let
\(A_0\) be their complement.
We need the following local identifiability bound. Hellinger distance
controls the mass outside the true-atom neighborhoods, as well as the first
two local moments of the perturbation inside each neighborhood. For
\[
 m_i=G(A_i),\qquad
 b_i=\int_{A_i}(\theta-\lambda_i)\,dG(\theta),\qquad
 s_i=\int_{A_i}(\theta-\lambda_i)^2\,dG(\theta),
\]
the Poisson analogue of the finite-mixture local-geometry bound is
\begin{equation}
\label{eq:poisson-local-geometry}
 G(A_0)+\sum_i\{|m_i-\pi_i|+|b_i|+s_i\}\le C\rho_G.
\end{equation}
We first prove the corresponding linear statement. Let \(\mu\) be a finite
nonnegative measure supported on \(A_0\), and let
\(\eta_i,\beta_i\in\mathbb R\) and \(\sigma_i\ge0\). The linearized density
perturbation is
\[
 p_\mu+\sum_i\{\eta_i p_{\lambda_i}
 +\beta_i\partial_\lambda p_{\lambda_i}
 +\tfrac12\sigma_i\partial_\lambda^2p_{\lambda_i}\},
\]
Its \(\ell^1\)-norm is bounded below by a fixed multiple of
\[
 \mu(A_0)+\sum_i(|\eta_i|+|\beta_i|+\sigma_i).
\]
Suppose this lower bound fails. Normalize
\(\mu(A_0)+\sum_i(|\eta_i|+|\beta_i|+\sigma_i)\) to one, and choose a
sequence whose linearized expression
tends to zero in \(\ell^1\). The coefficient vectors have a convergent
subsequence. The corresponding outer measures \(\mu_n\) are also tight.
Indeed, if \(v_n\)
denotes the finite linear combination of the fixed kernel derivatives, then
for every integer \(K\),
\[
 \sum_{k>K}p_{\mu_n}(k)
 \le \|p_{\mu_n}+v_n\|_1+\sum_{k>K}|v_n(k)|.
\]
The second term tends to zero uniformly in \(n\) as \(K\to\infty\).
Moreover,
\[
 \mu_n([R,\infty))\inf_{\theta\ge R}
       \mathbb P_\theta\{X>K\}
 \le \sum_{k>K}p_{\mu_n}(k),
\]
and the infimum tends to one as \(R\to\infty\), for fixed \(K\).
Thus a further subsequence satisfies \(\mu_n\Rightarrow\mu\) for a finite
measure \(\mu\) on \(A_0\). Passing to probability generating functions
gives the Laplace transform of
\[
 \mu+\sum_i\{\eta_i\delta_{\lambda_i}
 -\beta_i\delta'_{\lambda_i}
 +\tfrac12\sigma_i\delta''_{\lambda_i}\}
\]
identically zero. Uniqueness of Laplace transforms first eliminates \(\mu\)
away from the \(\lambda_i\)'s, and linear independence of the remaining
point masses and their derivatives eliminates every coefficient. This
contradicts the normalization.

To pass from linear identifiability to the nonlinear mixture bound, we use
the uniform Taylor remainder
\[
 \|p_\theta-p_{\lambda_i}-(\theta-\lambda_i)p'_{\lambda_i}
       -\tfrac12(\theta-\lambda_i)^2p''_{\lambda_i}\|_1
 \le C\,\operatorname{diam}(A_i)(\theta-\lambda_i)^2.
\]
This follows from
\(\partial_\theta^jp_\theta(k)=\sum_{r=0}^j(-1)^{j-r}
\binom jr p_\theta(k-r)\), with \(p_\theta(m)=0\) for \(m<0\), which
also gives a uniform \(\ell^1\)-bound on the third derivative. Shrinking
the fixed intervals \(A_i\) therefore proves the lower bound in total
variation.
Since total variation is at most twice Hellinger distance,
\eqref{eq:poisson-local-geometry} follows.

Let \(a=p_{G|_{A_0^c}}\) and \(s=\sqrt{a/p_0}\). Taylor expansion in the
fixed intervals and \eqref{eq:poisson-local-geometry} give
\begin{equation}
\label{eq:poisson-inner-outer-size}
 \|s-1\|_2\le C\rho_G,
 \qquad W:=G(A_0)\le C\rho_G.
\end{equation}
Fix this inner subdensity \(a\). We next bracket the increment
created by adding one outer component. Write \(q_u=p_{u^2}\), where
\(0\le u\le\sqrt M\), put \(b=wq_u\), and set
\(\nu(k)=b(k)^2/\{a(k)+b(k)\}\). After decreasing \(\rho_0\), the
local-geometry bound and \(\lambda_*>0\) ensure that the inner submixture
assigns a fixed positive mass to an interval
\([\ell,U]\subset(0,\infty)\). The ratios of consecutive moments of the
finite measure \(e^{-\theta}G|_{A_0^c}(d\theta)\) are nondecreasing, and
hence, using
\(a(k)=k!^{-1}\int e^{-\theta}\theta^k\,dG|_{A_0^c}(\theta)\),
\[
 \frac{a(k+1)}{a(k)}\ge\frac{c}{k+1},
 \qquad
 \frac{q_u(k+1)}{q_u(k)}=\frac{u^2}{k+1}\quad(u>0).
\]
It follows, by cross-multiplication, that, for \(u>0\),
\[
 (k+1)\nu(k+1)\le C(1+M)^2\nu(k).
\]
Put \(A=C(1+M)^2\). The recursion implies, for every \(t\ge0\),
\[
 \frac{\sum_ke^{tk}\nu(k)}{\sum_k\nu(k)}
 \le \exp\{A(e^t-1)\}.
\]
Indeed, the logarithmic derivative of the numerator is at most \(Ae^t\).
Taking \(t=c/A\) gives uniform exponential moments on count scale
\((1+M)^2\). For \(u\ge1/2\), differentiation of
\(p_{u^2}\) and Lemma~\ref{lem:contaminant-calculus} give
\[
 \|\partial_u h_{a,w,q_u}\|_2
 \le C(1+M)^2e_{a,w,q_u}.
\]
The interval \(0\le u\le1\) is handled directly. Indeed,
\(a(k)\ge c\ell^k/k!\), while termwise differentiation of
\(p_{u^2}(k)=e^{-u^2}u^{2k}/k!\) gives
\[
 \sup_{0\le u\le1}|\partial_u^jp_{u^2}(k)|
 \le C(k+1)/k!,\qquad j=0,1.
\]
These envelopes are square summable after division by \(a(k)\), and, for
the mixture weights \(0<w\le1\) considered here, the zeroth count
coordinate gives \(e_{a,w,q_u}\ge cw\). Hence the same
derivative bound holds on \([0,1]\), with a fixed constant.

On intervals of length \(c(1+M)^{-2}\) in \([1/2,\sqrt M]\), the exact
ratio
\[
 \frac{p_{(u+v)^2}(k)}{p_{u^2}(k)}
 =e^{-2uv-v^2}(1+v/u)^{2k}
 \le e^{C+Ck/(1+M)^2}
\]
and the exponential-moment bound for the count variable obtained above
control the pointwise suprema of the parameter derivatives. On \([0,1]\),
the power-series envelopes do the
same. The weight coordinate is model-free: Lemma
~\ref{lem:contaminant-calculus} shows that the logarithmic derivative of
the increment norm lies in \([1/2,1]\). Repeating the parameter-box and
weight-level construction in \eqref{eq:gaussian-one-contaminant-bracket}
therefore gives, uniformly in the inner subdensity,
\begin{equation}
\label{eq:poisson-one-contaminant-bracket}
 \log N_{[]}\left(
 \epsilon,
 \{h_{a,w,q_u}:0\le u\le\sqrt M,\ 0<w\le1,\
       e_{a,w,q_u}\le r\},
 L^2(p_0)\right)
 \le C\{\log M+\log_+(r/\epsilon)\}.
\end{equation}

The one-contaminant bound treats the inner subdensity as fixed. To count the
full mixture class, we must also bracket the inner atoms, all of which lie in
the fixed neighborhoods \(A_i\). We verify the needed entropy bound because
the Poisson family is not a translation family. Let
\(\mathcal D_D^{\rm in}\) be the normalized
nonzero Hellinger secants
\[
 \mathcal D_D^{\rm in}
 =\left\{
   \frac{\sqrt{p_H/p_0}-1}
        {\|\sqrt{p_H/p_0}-1\|_2}:
   0<\|\sqrt{p_H/p_0}-1\|_2\le\rho_0,
   \ |\supp(H)|\le D,\quad
   \supp(H)\subset\bigcup_iA_i
 \right\}.
\]
The proof of \cite[Theorem~3.1]{gassiatVanHandel2014} approximates precisely
such normalized secants in the Banach lattice \(L^2(p_0)\). As explained in
their Remark~3.6, translation structure is needed for their local-geometry
theorem, while the entropy argument itself applies to smooth parametric
kernels. Equation~\eqref{eq:poisson-local-geometry} supplies the required
local geometry. To verify smoothness, use the mean parameter and define, on
the fixed union of the \(A_i\)'s,
\[
 H_j(k)=\sup_\theta
       \frac{|\partial_\theta^jp_\theta(k)|}{p_0(k)},
 \qquad 0\le j\le3.
\]
The finite-difference formula for \(\partial_\theta^jp_\theta\), valid also
at \(\theta=0\), reduces each numerator to finitely many shifted Poisson
masses. If the inner intervals lie in \([0,B]\), then, outside a fixed
finite set of counts, their pointwise suprema are bounded by a polynomial
in \(k\) times \(p_B(k)\). Since
\(p_0\ge\pi_*p_{\lambda_*}\) for some \(\lambda_*>0\), the identity
\[
 \sum_{k\ge0}(1+k)^m
 \frac{p_B(k)^s}{p_{\lambda_*}(k)^{s-1}}<\infty,
 \qquad m<\infty,\quad s\in\{2,4\},
\]
yields
\[
 H_0,H_1,H_2\in\ell^4(p_0),
 \qquad H_3\in\ell^2(p_0).
\]
The pointwise Taylor bound and \eqref{eq:poisson-local-geometry} give the
analogue of \cite[Lemma~3.13]{gassiatVanHandel2014}. Specifically, the
density ratios normalized by their \(\ell^1\)-distance from \(p_0\) have
the \(\ell^4(p_0)\) envelope
\[
 \frac{|p_H/p_0-1|}{\|p_H-p_0\|_1}
 \le C(H_0+H_1+H_2).
\]
The algebraic comparison between normalized Hellinger and chi-square
secants in their Lemma~3.15 is measure-free. Their Lemma~3.16 treats atoms
close to a true mean by a second-order Taylor expansion and retains atoms of
small weight as individual kernel terms; the displayed \(H_j\) bounds
control the resulting remainder in \(\ell^2(p_0)\). Finally, pad the
candidate to \(D\vee K\) atoms and apply the coefficient and location grids
in the proof of their Theorem~3.1. The envelopes
\(H_0,H_1,H_2\in\ell^4(p_0)\), together with
\(H_3\in\ell^2(p_0)\), justify each Taylor remainder in counting measure,
and taking the pointwise
minimum and maximum over every parameter cell produces ordered brackets
exactly as in their proof. Those grids have at most \((CD/\xi)^{CD}\)
elements, and hence
\[
 \log N_{[]}\bigl(\xi,\mathcal D_D^{\rm in},L^2(p_0)\bigr)
 \le CD\log(CD/\xi).
\]
The general slicing argument in
\cite[Theorem~2.2]{gassiatVanHandel2014} converts this normalized-secant
bound into brackets of width \(\epsilon\) for the inner Hellinger ball of
radius \(r\), at entropy \(CD\log(CDr/\epsilon)\). Finally, write an inner
subdensity as \(a=m\bar a\). Then
\(|\sqrt m-1|\le\|\sqrt{a/p_0}-1\|_2\) and
\(\|\sqrt{\bar a/p_0}-1\|_2\le Cr\); bracketing \(\sqrt m\) separately
adds one scalar coordinate. This supplies the inner-submixture brackets
needed for the shell construction.

We now combine an inner bracket with at most \(D\) outer-component brackets,
using the composition map from Lemma~\ref{lem:contaminant-calculus}. On a shell
\(r/2<\rho_G\le r\), positivity and
\eqref{eq:poisson-inner-outer-size} give
\(\|h_{a,w_j,q_{u_j}}\|_2\le Cr\). Put
\(\zeta=c\min\{\eta,r\}\), bracket the inner square root and each outer
increment at width \(\zeta/(CD)\), and choose an actual baseline
representative from each nonempty inner bracket. Widen every component bracket
by the inner bracket width and compose through \(\Phi\). Since
\(\partial_s\Phi\le\sqrt{D+1}\) and
\(\partial_{h_j}\Phi\le1\), the resulting width is at most
\[
 \sqrt{D+1}\frac{\zeta}{CD}
 +D\left(\frac{\zeta}{CD}+2\frac{\zeta}{CD}\right)
 \le\zeta
\]
after increasing \(C\). Equations
\eqref{eq:poisson-one-contaminant-bracket} and the inner entropy bound give,
for the number \(N_r\) of brackets covering this shell,
\begin{equation}
\label{eq:poisson-shell-entropy}
 \log N_r\le CD\{\log M+\log D+\log_+(r/\eta)\}.
\end{equation}
For the innermost ball, Taylor expansion and
\eqref{eq:poisson-local-geometry} give a fixed \(F_0\in\ell^2(p_0)\) such
that the inner density ratio differs from one by at most \(C\rho_GF_0\).
The outer ratio is bounded by
\(C\rho_G\sup_{0\le\theta\le M}p_\theta/p_0\). Since
\(p_0\ge\pi_*p_{\lambda_*}\) and
\(\sup_{\theta\le M}p_\theta(k)=p_{k\wedge M}(k)\), Stirling's formula for
\(k\le M\) and the tilted-Poisson identity for \(k>M\) give a common
envelope of \(\ell^2(p_0)\)-norm at most \(e^{C(1+M^2)}\). Thus the ball
\(\rho_G\le\eta e^{-C(1+M^2)}\) fits in one bracket. Summing the dyadic shell
collections adds only \(C\{\log M+\log_+(\delta/\eta)\}\) to
\eqref{eq:poisson-shell-entropy}; \(D\le M^{C_0}\) absorbs \(D\log D\).

For \(\rho_G\ge\rho_0\), global \(L^1\) density brackets suffice. Use the
coordinate \(u=\sqrt\theta\), for which
\(\|\partial_up_{u^2}\|_1\le2\). For every interval
\(I\subset[0,\sqrt M]\), the pointwise endpoints
\[
 \ell_I(k)=\inf_{u\in I}p_{u^2}(k),
 \qquad u_I(k)=\sup_{u\in I}p_{u^2}(k)
\]
form a nonnegative order bracket. The fundamental theorem of calculus gives
\[
 \sum_k\{u_I(k)-\ell_I(k)\}
 \le\int_I\|\partial_up_{u^2}\|_1\,du\le2|I|.
\]
Partitioning \([0,\sqrt M]\) at mesh \(c\eta^2\) and quantizing the
weights at mesh \(c\eta^2/D\) produces mixture brackets of \(L^1\)-width
at most \(C\eta^2\). The number of location and weight choices is bounded
by \(\exp\{CD(\log M+\log(1/\eta))\}\). Thus
\[
 \log N_{[]}\bigl(\eta^2,
 \{p_G:|\supp(G)|\le D,\ \supp(G)\subset[0,M]\},L^1\bigr)
 \le CD\{\log M+\log(1/\eta)\}.
\]
Taking square roots and using \(\delta\ge\rho_0\) completes the proof.
\end{proof}

We now apply the localized likelihood inequality to the entropy bound. The
result controls all sparse Poisson competitors simultaneously and will be
compared with Proposition~\ref{prop:poisson-finite-constructive}.

\begin{proposition}
\label{prop:poisson-finite-upper}
Fix \(C_0<\infty\). Let \(D_n\) be positive integers satisfying
\(D_n\le(\log n)^{C_0}\) and \(D_n\log\log n=o(n)\). Then
\[
    \mathbb P\left\{
        \sup_{|\supp(G)|\le D_n}L_n^{\mathrm{Pois}}(G;G_0)
        \le C_{G_0,C_0}D_n\log\log n
    \right\}\to1.
\]
\end{proposition}

\begin{proof}
Let \(M_n=\max_iX_i\), and set \(m_n=\lfloor2\log n/\log\log n\rfloor\). Since \(G_0\) is fixed and finite, a Chernoff bound gives \(\mathbb P\{M_n\le m_n\}\to1\). By Lemma~\ref{lem:poisson-support}, on this event all \(D_n\)-atomic competitors may be restricted to \([0,m_n]\).

Let \(A_n=C_{G_0,C_0}D_n\log\log n\). Since
\(D_n\le m_n^{C_0+1}\) eventually and
\(\log m_n\le\log\log n+O(1)\), Lemma
~\ref{lem:poisson-finite-local}, with exponent \(C_0+1\), gives
\[
    \log N_{[]}
    \bigl(
        \eta,
        \mathcal F_{D_n}^{\mathrm{Pois},G_0}(m_n,\delta),
        L^2(p_0)
    \bigr)
    \le
    A_n(1+\log(\delta/\eta)).
\]
The localized bracketing integral is therefore \(O(\delta\sqrt{A_n})\), and Lemma~\ref{lem:wong-shen-local} proves the proposition.
\end{proof}

\begin{proof}[Proof of Theorem~\ref{thm:poisson-finite-main}]
Apply Proposition~\ref{prop:poisson-finite-constructive} with any fixed \(\gamma\in(0,1)\). Since \(J\asymp\sqrt{\log n/\log\log n}\), the unrestricted likelihood is at least \(c\sqrt{\log n/\log\log n}\) with probability tending to one. Proposition~\ref{prop:poisson-finite-upper} bounds every \(D_n\)-atomic competitor by \(C D_n\log\log n\) with high probability. Taking
\[
    D_n=
    \left\lfloor
        c_{G_0}\frac{\sqrt{\log n}}{(\log\log n)^{3/2}}
    \right\rfloor
\]
with \(c_{G_0}\) small enough, no \(D_n\)-atomic mixing distribution can maximize on the intersection of the two high-probability events. This proves the theorem.
\end{proof}

\small
\bibliographystyle{alpha}
\bibliography{writing_examples/bib,references}

\end{document}
